%====================================================
% Geometry of Sameness — Draft
%====================================================
\documentclass[11pt]{article}
\usepackage{amsmath}
% ---------- Encoding & Fonts ----------
\usepackage[T1]{fontenc}
\usepackage[utf8]{inputenc}
\usepackage{microtype}
\usepackage{lmodern}
\usepackage{newunicodechar}
\newunicodechar{→}{\textrightarrow{}}
\newunicodechar{∼}{\textasciitilde{}}
\newunicodechar{≤}{\ensuremath{\leq}}
\newunicodechar{−}{\ensuremath{-}}
\newunicodechar{≈}{\ensuremath{\approx}}

% ---------- Page Layout ----------
\usepackage[margin=1in]{geometry}
\usepackage{setspace}
\setstretch{1.08}

% ---------- Math & Symbols ----------
\usepackage{amsmath,amssymb,amsfonts,amsthm,bm}
\usepackage{mathtools}
\usepackage{enumitem}
\usepackage{tabularx}
\usepackage{siunitx}
\sisetup{detect-all=true}
\DeclareMathOperator*{\argmin}{arg\,min}
\DeclareMathOperator*{\argmax}{arg\,max}
\newcommand{\TSR}{\mathrm{TSR}}
\newcommand{\E}{\mathbb{E}}
\newcommand{\Prob}{\mathbb{P}}
\newcommand{\Var}{\mathrm{Var}}
\newcommand{\norm}[1]{\left\lVert #1 \right\rVert}
\newcommand{\RR}{\mathbb{R}}
\newcommand{\MSR}{\mathsf{MSR}}
% ---------- Graphics & Tables ----------
\usepackage{graphicx}
\usepackage{booktabs}
\usepackage{multirow}
\usepackage{adjustbox}

% ---------- Framed Boxes ----------
\usepackage{mdframed}
\usepackage{xcolor}

% ---------- Algorithms ----------
\usepackage{algorithm}
\usepackage{algpseudocode}
\usepackage{float}

\makeatletter
\renewcommand{\fps@algorithm}{htbp}
\setcounter{totalnumber}{10}
\setcounter{topnumber}{10}
\setcounter{bottomnumber}{10}
\renewcommand{\topfraction}{0.9}
\renewcommand{\bottomfraction}{0.9}
\renewcommand{\textfraction}{0.1}
\renewcommand{\floatpagefraction}{0.8}
\makeatother
\setlength{\intextsep}{10pt plus 2pt minus 2pt}

% ---------- Links & Clever Refs ----------
\usepackage[hidelinks]{hyperref}
\usepackage[capitalise,nameinlink]{cleveref}

% ---------- Quotes & Lists ----------
\usepackage{csquotes}

\setlist[itemize]{leftmargin=*, itemsep=2pt, topsep=2pt}
\setlist[enumerate]{leftmargin=*, itemsep=2pt, topsep=2pt}

% ---------- Theorem Envs ----------
\newtheorem{definition}{Definition}
\newtheorem{theorem}{Theorem}
\newtheorem{lemma}{Lemma}
\newtheorem{proposition}{Proposition}
\newtheorem{corollary}{Corollary}
\newtheorem{assumption}{Assumption}

% ---------- Convenience macros ----------
\newcommand{\SamTrans}{\mathsf{SamTrans}}
\newcommand{\SamGeom}{\mathsf{SamGeom}}

% ---------- Title Block ----------
\title{\textbf{The Geometry of Sameness:}\\[0.25ex]
\textbf{An \boldsymbol{$\varepsilon$}-Equivalence of Translation and Distance}\\[0.5ex]
\large Categories of Realizations, Error-Budget Transfer, and a Hidden Variable}
\author{Bee Rosa Davis\\
\small NASA Mission Systems Engineer\\
\small \texttt{bee\_davis@alumni.brown.edu}}
\date{}

\begin{document}
\maketitle
\vspace{-1.5ex}

\begin{abstract}
At an intuitive level, the problems we study are like trying to understand the Earth from maps. 
The underlying \emph{semantic sameness structure} $S$ is the globe: a hidden world of entities and 
their relationships. One engineering tradition starts from the globe (a manifold): learn a space 
where geodesic distance reflects semantic change and read decisions off geodesic gaps—this is the 
geometry-first view used by systems like HERALD and VIDAR. A second tradition starts from overlapping 
flat maps: heterogeneous observation spaces with translators between them, stitched together just 
enough to reconcile telemetry or transactions—this is the translation-first view used by systems like 
AMC and PRISM. Our central claim is that, whenever the “maps” are smooth and consistent enough to be 
sewn into a globe, these are not competing descriptions but dual realizations of the same $S$: tools 
and guarantees can be moved across the seam.

Formally, for a fixed semantic sameness structure $S$ we define two categories of realizations: \label{sec:categories} 
$\SamTrans(S)$, whose objects are \emph{translation-based} systems on heterogeneous observation spaces 
with bounded translator drift, and $\SamGeom(S)$, whose objects are \emph{manifold-based} systems with 
Riemannian metrics, path families, and configuration margins. We introduce \emph{$\varepsilon$-equivalence 
of categories} tailored to detection, together with explicit functors 
$F_S : \SamTrans^0(S) \to \SamGeom^0(S)$ and $G_S : \SamGeom^0(S) \to \SamTrans^0(S)$ that act like 
“stitching the maps into a globe” and “unfolding the globe back into maps” on well-behaved subcategories 
$\SamTrans^0(S), \SamGeom^0(S)$ that satisfy a \emph{smooth chartability} condition. Smooth chartability 
is the critical, falsifiable hypothesis of the framework: it is exactly the requirement that the local 
translators around $S$ can be approximated by a compatible atlas whose transition maps behave like 
local diffeomorphisms, so that a globe $M$ actually exists.

Our main result is an \emph{Error Budget Transfer Theorem}. Each realization (translator or manifold) 
comes with a four-term detection error budget 
$(E_{\mathrm{geom}}, E_{\mathrm{link}}, \xi, \zeta)$ for geometry, linkage, calibration, and abstention, 
together with an independence slack $\delta_{\mathrm{indep}}$ and a finite-variance Cantelli term 
$B_{\mathrm{Cantelli}}$ linking statistic separation to posterior correctness. We show that for any 
$T \in \SamTrans^0(S)$ and its geometric realization $M = F_S(T)$,
\[
E_{\mathrm{tot}}^{\mathrm{TSP}}(T)
\;=\;
1 - (1 - E_{\mathrm{geom}}^{\mathbf{T}})(1 - E_{\mathrm{link}}^{\mathbf{T}})
      (1 - \xi^{\mathbf{T}})(1 - \zeta^{\mathbf{T}})
      + \delta_{\mathrm{indep}}^{\mathbf{T}}
\]
and
\[
E_{\mathrm{tot}}^{\mathrm{MSP}}(M)
\;=\;
1 - (1 - E_{\mathrm{geom}}^{\mathbf{M}})(1 - E_{\mathrm{link}}^{\mathbf{M}})
      (1 - \xi^{\mathbf{M}})(1 - \zeta^{\mathbf{M}})
      + \delta_{\mathrm{indep}}^{\mathbf{M}}
\]
differ by at most an explicit, first-order slack 
$|E_{\mathrm{tot}}^{\mathrm{TSP}}(T) - E_{\mathrm{tot}}^{\mathrm{MSP}}(M)| \le C_F(\varepsilon_{\mathrm{trans}}, \varepsilon_{\mathrm{dist}}, \delta_{\mathrm{chart}})$, 
and the associated Cantelli bounds on high-risk alerts satisfy an analogous relation 
$|B_{\mathrm{Cantelli}}(T) - B_{\mathrm{Cantelli}}(M)| \le C_F'(\cdots)$. A symmetric statement holds 
for $M \in \SamGeom^0(S)$ and $T = G_S(M)$. In other words, once $S$ passes the smooth chartability 
test, \emph{translator-based and manifold-based realizations of $S$ share the same detection guarantees 
up to controlled slack}: the choice of “flat maps vs.\ globe” becomes an engineering decision, not a 
conceptual one.

Practically, this $\varepsilon$-equivalence has two consequences. First, it lifts theorems across the 
bridge: guarantees proved in a geometry-first system (e.g., Davis-style path-based bounds for HERALD 
or VIDAR) can be pulled back to translation-first architectures like AMC via $G_S$, and vice versa. 
Second, it lifts tools: audits developed for translator graphs (e.g., AMC’s spectral stability checks) 
become curvature or chart-overlap diagnostics on manifolds via $F_S$. We instantiate this framework 
on four systems—AMC and PRISM as translation-first, HERALD and VIDAR as geometry-first—to illustrate 
how chartability can be audited empirically and how error budgets and techniques transfer in practice.
\end{abstract}


%====================================================
\section{Introduction}

Many modern systems need to decide whether two heterogeneous observations refer to the same
underlying entity or event: is this telemetry packet semantically equivalent to that one after
compression and format switching?  Does this payment record match that ledger entry across rails?
Do these viral variants represent the same antigenic state?  Does this video clip belong to the
same physical person as that enrollment image?

Two families of architectures answer these questions in different languages:

\begin{itemize}
  \item \textbf{Translator-based systems} (e.g., AMC-style adaptive format pipelines, multi-rail
  reconciliation) build a graph of learned translators between embedding formats or modalities and
  assert sameness when translated representations agree within a drift budget.
  \item \textbf{Manifold-based systems} (e.g., Davis-style HERALD and VIDAR) learn or inherit a
  Riemannian manifold on which geodesic distance along identity-preserving paths encodes semantic
  change, and assert sameness when points remain within a configuration margin in that geometry.
\end{itemize}

Informally, one view speaks in \emph{translations} and drift; the other in \emph{distances} and
curvature.  Both solve the same kind of problem, but their guarantees, diagnostics, and engineering
tools are typically developed in isolation.

This work starts from a simple question:

\begin{quote}
\emph{For a fixed notion of semantic sameness, are translator-based and manifold-based detection
systems two views of the same underlying problem?  If so, can guarantees and tools proven in one
view be transported to the other without re-deriving them from scratch?}
\end{quote}

We answer this in the affirmative under explicit regularity assumptions.  The key move is to make
the \emph{problem} itself explicit.

\subsection{Semantic sameness structures and realizations}\label{sec:sameness}

We model the underlying task as a \emph{semantic sameness structure} $S$:
latent identities, observation spaces, rendering maps, and a sameness relation on observations.  For
a fixed $S$, engineers can build many concrete systems.

\begin{itemize}
  \item $\SamTrans(S)$: a category whose objects are \emph{translator-based realizations} of $S$
  (observation encoders, format/rail translators, drift budgets, and detection logic in translator
  space), and whose morphisms are reparameterizations preserving translators up to controlled
  error.
  \item $\SamGeom(S)$: a category whose objects are \emph{manifold-based realizations} of $S$
  (Riemannian manifolds, chart atlases, path families, and Davis-style detection logic), and whose
  morphisms are quasi-isometries respecting configuration structure and charts.
\end{itemize}

Crucially, these are \emph{two} design spaces for the same $S$: translator systems and manifold
systems are not being unified with each other directly, but as two families of realizations of the
same hidden variable.

On well-behaved subcategories $\SamTrans^0(S) \subset \SamTrans(S)$ and
$\SamGeom^0(S) \subset \SamGeom(S)$, defined by smooth chartability, bounded distortion, and
non-vacuous configuration margins, we construct functors
\[
F_S : \SamTrans^0(S) \to \SamGeom^0(S),
\qquad
G_S : \SamGeom^0(S) \to \SamTrans^0(S),
\]
that turn translators into manifolds and manifolds into translator graphs.  These functors are not
magic; they succeed precisely when the chartability and distortion assumptions are satisfied.

\subsection{Main contributions}

Our contributions are organized around three load-bearing claims and their practical implications.

\paragraph{(0) Language: $\varepsilon$-equivalence of realization categories.}
We formalize translator and manifold realizations of a fixed semantic sameness structure $S$ as
categories $\SamTrans(S)$ and $\SamGeom(S)$ and introduce an \emph{$\varepsilon$-equivalence} notion
tailored to detection.  Rather than claiming a strict equivalence of categories, we construct
functors $F_S$ and $G_S$ and show that
$G_S \circ F_S \simeq \mathrm{Id}_{\SamTrans^0(S)}$ and
$F_S \circ G_S \simeq \mathrm{Id}_{\SamGeom^0(S)}$ \emph{up to quasi-isometry and bounded morphism
error}, with all slack terms explicit.  Categorical language is not decorative here: it is the
minimal formalism needed to express how \emph{error budgets} transform under representation change.

\paragraph{(1) Error Budget Transfer Theorem (main result).}
Our central technical result shows that detection guarantees transfer across the two
realizations. Conceptually, we prove that the correctness bound for a translator system
($T$) is \(\varepsilon\)-equivalent to the correctness bound for the manifold system ($M$)
that realizes the same problem $S$:
%
% This is the "sexy" equation
\begin{equation*}
\mathrm{LB}(T) \cong_{\varepsilon} \mathrm{LB}(F_S(T))
\end{equation*}
%
This main result, the Error Budget Transfer Theorem, is stated formally below. It shows
that the full, multi-part error budgets for both systems are equivalent up to explicit,
first-order slack terms.

\vspace{1ex}
% This is the formal, boxed theorem
\begin{mdframed}[backgroundcolor=gray!10, linewidth=0.4pt]
\noindent\textbf{Error Budget Transfer Theorem (informal preview).}
\emph{Fix a semantic sameness structure $S$ and consider well-behaved realizations
$T \in \SamTrans^0(S)$ and $M \in \SamGeom^0(S)$.}

\emph{Translator-side bound.} The translator realization $T$ admits a Davis-style correctness
lower bound
\[
\mathrm{LB}_{\mathrm{T}}(T)
= B_{\mathrm{Cantelli}}(T)\,
\Bigl[\,(1-E_{\mathrm{geom}}^{\mathrm{T}})
(1-E_{\mathrm{link}}^{\mathrm{T}})
(1-\xi^{\mathrm{T}})
(1-\zeta^{\mathrm{T}})
- \delta_{\mathrm{indep}}^{\mathrm{T}}\Bigr]_+
- \varepsilon_{\mathrm{est}}^{\mathrm{T}}.
\]

\emph{Manifold-side bound.} The manifold realization $M$ admits an analogous bound
\[
\mathrm{LB}_{\mathrm{M}}(M)
= B_{\mathrm{Cantelli}}(M)\,
\Bigl[\,(1-E_{\mathrm{geom}}^{\mathrm{M}})
(1-E_{\mathrm{link}}^{\mathrm{M}})
(1-\xi^{\mathrm{M}})
(1-\zeta^{\mathrm{M}})
- \delta_{\mathrm{indep}}^{\mathrm{M}}\Bigr]_+
- \varepsilon_{\mathrm{est}}^{\mathrm{M}}.
\]

\emph{Transfer.} The functors $F_S$ and $G_S$ transport these guarantees:
\[
\bigl|\mathrm{LB}_{\mathrm{T}}(T) - \mathrm{LB}_{\mathrm{M}}(F_S(T))\bigr|
\;\le\;
C_F(\varepsilon_{\mathrm{trans}}, \varepsilon_{\mathrm{dist}}, \delta_{\mathrm{chart}}),
\]
\[
\bigl|\mathrm{LB}_{\mathrm{M}}(M) - \mathrm{LB}_{\mathrm{T}}(G_S(M))\bigr|
\;\le\;
C_G(\varepsilon_{\mathrm{trans}}, \varepsilon_{\mathrm{dist}}, \delta_{\mathrm{chart}}),
\]
\emph{where $C_F$ and $C_G$ are explicit first-order slack functions.}
\end{mdframed}
\vspace{1ex}

\paragraph{(2) Characterization via smooth chartability.}
The functor $F_S : \SamTrans^0(S) \to \SamGeom^0(S)$ is not automatic; it succeeds only on translator
systems with \emph{smooth chartability}.  We isolate this as the central, falsifiable hypothesis of
the framework: locally linear approximants to each translator must exist with bounded Jacobian and
small residual drift on in-regime data.  This assumption is testable via concrete audit protocols
(e.g., Jacobian-based linearization error histograms on AMC-style translators) and is empirically
supported in systems whose architectures already regularize for smoothness and stability.

\paragraph{(3) Constructive functors and $\varepsilon$-equivalence.}
Under smooth chartability and bounded distortion, we give explicit constructions of
$F_S$ (path-metric completion of a translator graph followed by a Riemannian smoothing) and
$G_S$ (chart-transition translators induced by overlapping coordinate maps).  We show that on
$\SamTrans^0(S)$ and $\SamGeom^0(S)$ these functors form an $\varepsilon$-equivalence of categories:
they are quasi-inverse up to explicitly bounded drift, curvature, and chart-overlap slack.  This
formalizes the intuition that translator and manifold views are dual coordinatizations of the same
semantic sameness problem $S$.

\paragraph{Implications: tool and theorem transfer.}
The practical consequence of this $\varepsilon$-equivalence is that \emph{tools and theorems} proven
in one view can now be reused in the other.

\begin{itemize}
  \item \emph{Theorems transfer.}  A guarantee proven in a manifold system (e.g., HERALD's
  Davis-style dominance lower bound, VIDAR's deepfake correctness bounds) can be pulled back via
  $G_S$ to a translator-based realization of the same $S$ (e.g., AMC-style telemetry translators),
  with only first-order changes in the error budget.
  \item \emph{Tools transfer.}  Engineering techniques developed for translator graphs (spectral
  stability audits, drift-triggered retraining, lineage-aware rollback) can be pushed forward via
  $F_S$ to become curvature and chart-stability diagnostics on manifolds, providing new audit
  tools for geometry-first systems.
\end{itemize}

Engineers starting from a single semantic sameness structure $S$ therefore have a principled choice
of working in translator or manifold coordinates, knowing that both design spaces are linked by
functors that preserve detection guarantees up to explicitly quantified slack.


\subsection{Running examples}

We ground the theory in four concrete systems, which we treat as realizations of different semantic
sameness structures $S$.

\paragraph{AMC (telemetry format switching).}
AMC~\cite{davis2025amc} is a telemetry-native embedding architecture for adaptive format switching 
and semantic translation in high-volume telemetry pipelines. It maintains a directed graph of 
learned format translators constrained by an $\varepsilon$-bounded semantic drift loss and uses 
spectral properties of a format-compatibility graph to select stable routes and fallbacks. In our 
language, AMC is a canonical object of $\SamTrans(S_{\mathrm{tele}})$ for a telemetry sameness 
structure $S_{\mathrm{tele}}$.

\paragraph{PRISM (payment reconciliation).}
PRISM~\cite{davis2025prism} performs multi-rail payment reconciliation with learned translators 
between heterogeneous transaction formats, semantic alignment constraints, and a central latent 
reconciliation space. It naturally lives in $\SamTrans(S_{\mathrm{pay}})$ but exposes an implicit 
manifold structure over reconciled transaction identities.

\paragraph{HERALD (viral antigenic drift).}
HERALD~\cite{davis2025herald}, introduced in prior work on Davis manifolds~\cite{davis2025manifold}, 
builds a pullback Riemannian manifold on viral sequences where geodesic distance along 
epitope-restricted mutation paths approximates antigenic distance; a PIT-fused drift statistic 
and Cantelli bounds map geometric drift to dominance risk with an explicit error budget. 
HERALD is a prototypical object of $\SamGeom(S_{\mathrm{antigen}})$.\,\textsuperscript{\!}\!

\paragraph{VIDAR (deepfake detection via identity trajectories).}
VIDAR treats face-recognition embeddings as points on the hypersphere, models short clips as
identity trajectories on an identity manifold, and builds RIM features plus auxiliary detectors as
inputs to a Davis-style scalar statistic and bound.\,\footnote{See Davis,
``The Davis Manifold: Geometry-First Detection with Compositional Error Budgets,'' for the VIDAR
instantiation.}  VIDAR is a canonical object of $\SamGeom(S_{\mathrm{id}})$.

Across these examples, the same pattern recurs: translator systems like AMC and PRISM live in
$\SamTrans(S)$; geometry-first systems like HERALD and VIDAR live in $\SamGeom(S)$; and our
functors $F_S$ and $G_S$ explain how to move guarantees and tools between these views.

\subsection{Scope and non-goals}

This manuscript is deliberately theoretical.  It:

\begin{itemize}
  \item defines semantic sameness structures $S$ and the categories $\SamTrans(S)$ and
  $\SamGeom(S)$;
  \item introduces smooth chartability and bounded-distortion path geometry as explicit, falsifiable
  assumptions defining subcategories $\SamTrans^0(S)$ and $\SamGeom^0(S)$;
  \item constructs $F_S$ and $G_S$ and proves an $\varepsilon$-equivalence of realization categories;
  \item proves the Error Budget Transfer Theorem (Theorem~6.3) showing that Davis-style correctness
  bounds transfer across these views with first-order slack; and
  \item sketches audit protocols for checking smooth chartability and error-budget components in
  systems like AMC, PRISM, HERALD, and VIDAR.
\end{itemize}

We do \emph{not} report new empirical performance, deployable calibration curves, or production
latencies; those belong in system-specific papers.  Our goal here is to expose the common hidden
variable $S$, make the representation change between translator and manifold views explicit and
functorial, and characterize when and how detection guarantees can be transported between them.

\subsection{Roadmap}

The paper develops in three acts.

\textbf{Act I (Framework).}  Section~2 formalizes semantic sameness structures $S$ and defines
translator and manifold realizations, together with Davis-style error budgets and the notion of
smooth chartability.  Section~3 introduces the categories $\SamTrans(S)$ and $\SamGeom(S)$,
constructs the functors $F_S$ and $G_S$ at a high level, and defines $\varepsilon$-equivalence of\label{def:eps-equivalence}
realization categories.

\textbf{Act II (Constructions and transfer).}  Section~4 builds $F_S$ concretely by endowing
translator graphs with a path metric and a smoothed Riemannian structure under smooth chartability,
and analyzes the resulting distortion and chart-overlap slack.  Section~5 constructs $G_S$ from
manifold charts and transition maps and analyzes its stability.  Section~6 proves the Error Budget
Transfer Theorem (Theorem~6.3) by tracking how geometric, linkage, calibration, and abstention
errors transform under $F_S$ and $G_S$.

\textbf{Act III (Case studies and outlook).}  Section~7 instantiates the framework on AMC, PRISM,
HERALD, and VIDAR, including concrete chartability audits and error-budget estimates.  Section~8
discusses limitations, open problems, and the broader “Davis universality” conjecture for
geometry-first detection in temporal semantic sameness problems.

\section{Semantic sameness structures and realizations}
\label{sec:sameness-structure}

Throughout this section we fix a \emph{semantic sameness problem} and make precise what it
means to ``realize'' that problem either as a translator graph or as a Riemannian manifold.
The hidden object is the sameness structure itself; translation-based and manifold-based
systems are two different realizations of the same underlying task.

We use boldface symbols such as $\mathbf{T}$ and $\mathbf{M}$ for concrete systems
(a particular TSR or MSR) and reserve plain letters for abstract spaces and maps.

\subsection{Semantic sameness structures}
\label{subsec:semantic-structure}

Intuitively, a semantic sameness structure records:
(i) a latent space of entities $u$ (viruses, identities, vehicles, \dots),
(ii) how each modality $i$ renders those entities as observations $x_i$,
(iii) when two observations in possibly different modalities ``are the same entity'', and
(iv) which latent trajectories count as benign identity-preserving evolution.

\begin{definition}[Semantic sameness structure]\label{def:sameness-structure}
\label{def:semantic-structure}
A \emph{semantic sameness structure} is a tuple
\[
S := \bigl( I,\; \mathcal{I},\; \{X_i\}_{i\in\mathcal{I}},\; \{\pi_i\}_{i\in\mathcal{I}},\; \approx,\;
\{\mathcal{P}_S(L)\}_{L>0} \bigr),
\]
where:
\begin{enumerate}[label=\textnormal{(\roman*)}, leftmargin=*]
\item $I$ is a (typically high-dimensional) set of \emph{latent entities} or states. A point
$u\in I$ represents the underlying ``object of interest'' (e.g., true antigenic state, true
human identity, true vehicle configuration).

\item $\mathcal{I}$ is a finite index set of \emph{modalities} (e.g., file formats, sensors,
assays, views).

\item For each $i\in\mathcal{I}$, $X_i$ is an \emph{observation space} for modality $i$.
We write
\[
X := \bigsqcup_{i\in\mathcal{I}} X_i
\]
for the disjoint union of all observations.

\item For each $i\in\mathcal{I}$, $\pi_i : I \to X_i$ is a (possibly partial) \emph{rendering
map} that produces an ideal observation of $u$ in modality $i$ when defined.\footnote{%
Formally, one may model $\pi_i$ as a Markov kernel from $I$ to $X_i$, capturing stochastic
observation noise; for notational simplicity we write $\pi_i$ as a map and treat noise
separately in the detector error budget.}
When $\pi_i(u)$ is undefined, modality $i$ simply does not observe $u$.

\item $\approx$ is a \emph{semantic sameness relation} on $X$, defined by
\[
x \approx x'
\quad\Longleftrightarrow\quad
\exists u\in I,\; \exists i,j\in\mathcal{I} \text{ such that }
x \text{ is an observation ``of'' } u \text{ in } X_i,\;
x' \text{ is an observation ``of'' } u \text{ in } X_j.
\]
In the noiseless idealization this reduces to:
$x\approx x'$ if there exist $u\in I$ and $(i,j)$ with $x=\pi_i(u)$ and $x'=\pi_j(u)$.
In practice we blur small observation noise and treat $x$ as an observation ``of'' $u$
whenever it lies in a pre-specified neighborhood of $\pi_i(u)$.

\item For each $L>0$, $\mathcal{P}_S(L)$ is a family of \emph{benign latent paths}
$\gamma : [0,1]\to I$ indexed by a \emph{semantic horizon} $L$.
Each $\gamma\in\mathcal{P}_S(L)$ represents an identity-preserving evolution of latent state
(e.g., realistic antigenic drift over one season, smooth identity motion over a short clip).
We require that $L\mapsto \mathcal{P}_S(L)$ is monotone:
if $L_1\le L_2$ then $\mathcal{P}_S(L_1)\subseteq \mathcal{P}_S(L_2)$.
\end{enumerate}
\end{definition}

\paragraph{Examples (semantic structures).}
\begin{itemize}[leftmargin=*]
\item \textbf{Format-heterogeneous telemetry (AMC~\cite{davis2025amc}).} $I$ is the space of 
underlying system states; each $X_i$ is a particular log or telemetry format; $\pi_i$ renders 
a state into that format; $\gamma\in\mathcal{P}_S(L)$ is a short-time state trajectory.

\item \textbf{Antigenic drift (HERALD).} $I$ are viral lineages; $X_{\mathrm{seq}}$ are amino-acid
sequences; $X_{\mathrm{neut}}$ are neutralization measurements; $\pi_{\mathrm{seq}}$ encodes the
canonical genome for a lineage; $\pi_{\mathrm{neut}}$ renders idealized assay readouts;
$\mathcal{P}_S(L)$ are mutation paths with at most $L$ epitope substitutions.

\item \textbf{Identity trajectories (VIDAR).} $I$ are person identities; $X_{\mathrm{video}}$
are short talking-head clips; $X_{\mathrm{audio}}$ are speech segments; each $\pi_i$ renders a
clip of the same person; $\mathcal{P}_S(L)$ are natural motion trajectories over windows of
bounded duration.
\end{itemize}

We will later view \emph{realizations} of a fixed $S$ as objects of two categories:
a translation-based category $\SamTrans(S)$ and a geometry-based category $\SamGeom(S)$.

\subsection{Translation-based realizations (TSRs)}
\label{subsec:tsr}

A translation-based realization starts from modality-specific feature spaces and implements
sameness comparisons via learned translators between those spaces.

\begin{definition}[Translation-based semantic realization (TSR)]
\label{def:tsr}
Let $S$ be a semantic sameness structure as in Definition~\ref{def:semantic-structure}.
A \emph{translation-based semantic realization} (TSR) of $S$ is a tuple
\[
\mathbf{T}
:=
\bigl(
\{V_i\}_{i\in\mathcal{I}},
\{\varphi_i\}_{i\in\mathcal{I}},
G_T,
\{T_{ij}\}_{(i,j)\in E_T},
\varepsilon_{\mathrm{trans}}(\cdot)
\bigr),
\]
where:
\begin{enumerate}[label=\textnormal{(\roman*)}, leftmargin=*]
\item For each $i\in\mathcal{I}$, $V_i$ is a finite-dimensional inner-product space
(feature space for modality $i$).

\item $\varphi_i : X_i\to V_i$ is a feature extractor (e.g., a neural encoder).
We denote ideal features by
\[
v_i(t) := \varphi_i\bigl(\pi_i(\gamma(t))\bigr)
\qquad
\text{for } \gamma\in\mathcal{P}_S(L),\; t\in[0,1],
\]
whenever $\pi_i(\gamma(t))$ is defined.

\item $G_T = (\mathcal{I}, E_T)$ is a directed graph whose vertices are modalities and whose
edges $i\to j$ indicate the availability of a learned \emph{translator}
$T_{ij} : V_i \to V_j$.

\item For any path $p = (i_0,\dots,i_k)$ in $G_T$ we write
\[
T_p := T_{i_{k-1} i_k} \circ \cdots \circ T_{i_0 i_1}
\quad\text{and}\quad |p|:=k.
\]

\item $\varepsilon_{\mathrm{trans}} : \mathbb{N} \to [0,\infty)$ is a \emph{translator drift profile}
such that, for any benign latent path $\gamma\in\mathcal{P}_S(L)$, any time $t\in[0,1]$,
any modalities $i,j$ with both $\pi_i(\gamma(t))$ and $\pi_j(\gamma(t))$ defined,
and any $G_T$-path $p: i\to j$ of length $k$, we have
\begin{equation}
\label{eq:tsr-drift}
\bigl\|
T_p\bigl( v_i(t) \bigr)
-
v_j(t)
\bigr\|
\;\le\;
\varepsilon_{\mathrm{trans}}(k).
\end{equation}
This inequality measures \emph{representation error} of the translators at a single latent
time, decoupled from any semantic drift along $\gamma$.
\end{enumerate}
We write $\TSR(S)$ for the collection of all TSRs that realize $S$.
\end{definition}

\paragraph{Examples (TSRs).}
\begin{itemize}[leftmargin=*]
\item \textbf{AMC-style telemetry~\cite{davis2025amc}.} $V_i$ are format-specific embedding 
spaces; $\varphi_i$ encode logs into these spaces; $T_{ij}$ are learned format translators; 
$\varepsilon_{\mathrm{trans}}$ bounds accumulated translator drift over paths in the format graph.

\item \textbf{PRISM-style multi-signal systems~\cite{davis2025prism}.} $V_i$ are detector-specific 
representations; $T_{ij}$ fuse or reconcile outputs between detectors; $\varepsilon_{\mathrm{trans}}$ 
captures how far translations can drift while still preserving semantic identity.
\end{itemize}

\subsection{Manifold-based realizations (MSRs)}
\label{subsec:msr}

A manifold-based realization represents semantic sameness by embedding latent entities into
a Riemannian manifold and requiring that feature-space coordinates act as approximate charts.

\begin{definition}[Manifold-based semantic realization (MSR)]
\label{def:msr}
Let $S$ be as in Definition~\ref{def:semantic-structure}.
A \emph{manifold-based semantic realization} (MSR) of $S$ is a tuple
\[
\mathbf{M}
:=
\bigl(
\{V_i\}_{i\in\mathcal{I}},
\{\varphi_i\}_{i\in\mathcal{I}},
(\mathcal{M}, g),
\{\psi_i\}_{i\in\mathcal{I}},
\varepsilon_{\mathrm{dist}}(\cdot)
\bigr),
\]
where:
\begin{enumerate}[label=\textnormal{(\roman*)}, leftmargin=*]
\item $\{V_i\}_{i\in\mathcal{I}}$ and $\{\varphi_i : X_i\to V_i\}_{i\in\mathcal{I}}$ are
feature spaces and encoders as in a TSR.

\item $(\mathcal{M}, g)$ is a $d$-dimensional smooth Riemannian manifold of \emph{semantic
states} with geodesic distance $d_g$.

\item For each $i\in\mathcal{I}$, $\psi_i : V_i \to \mathcal{M}$ is a smooth \emph{parameterization}
whose image $U_i := \psi_i(V_i)$ is an open subset of $\mathcal{M}$. We require that the
inverse maps
\[
\chi_i := \psi_i^{-1} : U_i \to V_i
\]
exist and are smooth; the family $\{(U_i,\chi_i)\}_{i\in\mathcal{I}}$ forms a (possibly
partial) chart atlas on the portion of $\mathcal{M}$ we actually use.

\item There exists a \emph{latent state map} $F : I \to \mathcal{M}$ such that, whenever
$\pi_i(u)$ is defined,
\[
F(u) = \psi_i\bigl( \varphi_i(\pi_i(u)) \bigr).
\]
This ensures that, on ideal data, all modalities agree on the same manifold point
for a given latent entity.

\item $\varepsilon_{\mathrm{dist}} : (0,\infty)\to [0,\infty)$ is a \emph{metric distortion
profile} such that, for any chart $(U_i,\chi_i)$, any two points $z_a,z_b\in U_i$ with
$d_g(z_a,z_b)>0$, and chart distance
\[
\delta_i(z_a,z_b)
:=
\bigl\| \chi_i(z_a) - \chi_i(z_b) \bigr\|_{V_i},
\]
we have
\begin{equation}
\label{eq:msr-distortion}
\left|
\frac{\delta_i(z_a,z_b)}{d_g(z_a,z_b)} - 1
\right|
\;\le\;
\varepsilon_{\mathrm{dist}}\bigl( d_g(z_a,z_b) \bigr).
\end{equation}
This $\varepsilon_{\mathrm{dist}}$ measures \emph{representation error} of the charts
$\chi_i$ relative to the underlying metric $g$.
\end{enumerate}
We write $\MSR(S)$ for the collection of all MSRs that realize $S$.
\end{definition}

\paragraph{Examples (MSRs).}
\begin{itemize}[leftmargin=*]
\item \textbf{HERALD-style geometry.} $(\mathcal{M}, g)$ is an antigenic manifold; $V_i$
is the latent space of a sequence encoder; $\psi_i$ is the learned map into $\mathcal{M}$;
$\varepsilon_{\mathrm{dist}}$ captures geodesic--Euclidean distortion along antigenic trajectories.

\item \textbf{VIDAR-style identity manifold.} $(\mathcal{M}, g)$ is the hypersphere
$\mathbb{S}^{d-1}$ with its standard metric; $V_{\mathrm{video}}$ is the pre-normalized
ArcFace embedding space; $\psi_{\mathrm{video}}$ is $\ell_2$-normalization; $\varepsilon_{\mathrm{dist}}$
is small in the small-angle regime.
\end{itemize}

\subsection{Semantic drift profile}
\label{subsec:semantic-drift}

Once an MSR $\mathbf{M}$ is fixed, latent benign paths in $\mathcal{P}_S(L)$ induce paths on
$\mathcal{M}$ via the latent map $F$.

\begin{definition}[Semantic drift profile]
\label{def:semantic-drift}
Let $S$ be a semantic sameness structure and $\mathbf{M}$ an MSR of $S$ as in
Definition~\ref{def:msr}. For $L>0$ and any $\gamma\in\mathcal{P}_S(L)$ define the induced
manifold path
\[
z_\gamma(t) := F\bigl(\gamma(t)\bigr) \in \mathcal{M}.
\]
The \emph{semantic drift profile} of $S$ under $\mathbf{M}$ is
\[
\sigma_S^{\mathbf{M}}(L)
:=
\sup_{\gamma\in\mathcal{P}_S(L)}\;
\sup_{0\le s\le t\le 1}
d_g\bigl(z_\gamma(s), z_\gamma(t)\bigr).
\]
Thus $\sigma_S^{\mathbf{M}}(L)$ records, in purely geometric terms, how far benign latent
paths are allowed to move in $\mathcal{M}$ over horizon $L$. It is a property of the pair
$(S,\mathbf{M})$, independent of translators or observation noise.
\end{definition}

\subsection{Chart overlap slack and translator asymmetry}
\label{subsec:chart-asymm}

When a single $S$ admits both a TSR $\mathbf{T}$ and an MSR $\mathbf{M}$, we will compare them
via two structural diagnostics: a chart-overlap slack that measures how well translators agree
with chart transitions, and a translator asymmetry that measures round-trip drift.

\begin{definition}[Chart overlap slack]
\label{def:chart-slack}
Let $S$ be fixed, $\mathbf{T}$ a TSR and $\mathbf{M}$ an MSR of $S$ sharing the same
feature spaces $\{V_i\}_{i\in\mathcal{I}}$ and encoders $\{\varphi_i\}_{i\in\mathcal{I}}$.
For each $i$, let $U_i:=\psi_i(V_i)\subseteq\mathcal{M}$ and $\chi_i:=\psi_i^{-1}:U_i\to V_i$.

For any overlap region $U_i\cap U_j\neq\varnothing$ and point $z\in U_i\cap U_j$,
the ideal coordinate change from modality $i$ to modality $j$ maps
\[
v_i := \chi_i(z)
\quad\mapsto\quad
v_j^\star := \chi_j(z).
\]
The translator $T_{ij}$ instead produces $T_{ij}(v_i)$. The \emph{chart overlap slack} is
\[
\delta_{\mathrm{chart}}(\mathbf{T},\mathbf{M})
:=
\sup_{(i,j)\in E_T}
\;\sup_{z\in U_i\cap U_j}
\bigl\|
T_{ij}\bigl(\chi_i(z)\bigr) - \chi_j(z)
\bigr\|.
\]
This quantity measures how far translator-based coordinate changes deviate from the
chart transitions implied by $\mathbf{M}$.
\end{definition}

\begin{definition}[Translator asymmetry]
\label{def:translator-asymmetry}
Fix a TSR $\mathbf{T}$ with graph $G_T=(\mathcal{I},E_T)$ and feature spaces $\{V_i\}$.
Let $L_0>0$ be a chosen operating horizon and define, for each ordered pair $(i,j)$ with
both $\pi_i$ and $\pi_j$ defined along $\mathcal{P}_S(L_0)$,
\[
\Omega_{ij}(\mathbf{T})
:=
\Bigl\{
v_i(t)
=
\varphi_i\bigl(\pi_i(\gamma(t))\bigr)
\;\Big|\;
\gamma\in\mathcal{P}_S(L_0),\; t\in[0,1],\;
\pi_j(\gamma(t))\ \text{defined}
\Bigr\}
\subseteq V_i.
\]
Assume that $G_T$ contains edges $i\to j$ and $j\to i$. The \emph{translator asymmetry}
of $\mathbf{T}$ is
\[
\delta_{\mathrm{asymm}}(\mathbf{T})
:=
\sup_{(i,j)\in E_T,\; (j,i)\in E_T}
\;\sup_{v\in \Omega_{ij}(\mathbf{T})}
\bigl\|
T_{ji}\bigl(T_{ij}(v)\bigr) - v
\bigr\|.
\]
This is a structural parameter: it measures worst-case round-trip drift on in-regime
features, not a probability. It will enter later as a diagnostic and, when needed,
as a slack term in error-transfer bounds.
\end{definition}

\subsection{Sameness detectors and error budgets}
\label{subsec:detector}

Finally, we formalize detectors that operate on a fixed realization $\mathbf{R}$ of $S$ and
attach a decomposed error budget. Here $\mathbf{R}$ may be either a TSR $\mathbf{T}$ or an
MSR $\mathbf{M}$.

\begin{definition}[Sameness detector and error budget]
\label{def:sameness-detector}
Let $S$ be a semantic sameness structure and $\mathbf{R}$ a realization of $S$ (TSR or MSR).
A \emph{sameness detector} on $\mathbf{R}$ is a tuple
\[
\mathbf{D_R}
:=
\bigl(
\Sigma_{\mathbf{R}},\,
A_{\mathbf{R}},\,
E_{\mathrm{geom}}^{\mathbf{R}},\,
E_{\mathrm{link}}^{\mathbf{R}},\,
\xi^{\mathbf{R}},\,
\zeta^{\mathbf{R}},\,
\delta_{\mathrm{indep}}^{\mathbf{R}},\,
\tau_{\mathrm{vac}}
\bigr),
\]
with components:
\begin{enumerate}[label=\textnormal{(\roman*)}, leftmargin=*]
\item $\Sigma_{\mathbf{R}}$ is a scalar \emph{detection or similarity statistic} built from the
representation induced by $\mathbf{R}$:
\begin{itemize}
\item for a TSR $\mathbf{T}$, $\Sigma_{\mathbf{T}}$ acts on feature-space objects
(e.g., paths or pairs in $\bigsqcup_i V_i$);
\item for an MSR $\mathbf{M}$, $\Sigma_{\mathbf{M}}$ acts on points or paths in $\mathcal{M}$.
\end{itemize}

\item $A_{\mathbf{R}}$ is a decision rule mapping $(\Sigma_{\mathbf{R}}, \text{context})$ to a
finite action set such as $\{\mathsf{same},\mathsf{changed},\mathsf{ambiguous},\mathsf{abstain}\}$
or $\{\mathsf{info},\mathsf{warn},\mathsf{high\_risk},\mathsf{insufficient\_signal}\}$.

\item $E_{\mathrm{geom}}^{\mathbf{R}}$ is the \emph{geometric / representation error}:
the probability (on in-regime, non-abstaining inputs) that the structural assumptions of
$\mathbf{R}$ fail on ideal data:
\begin{itemize}
\item for a TSR $\mathbf{T}$, $E_{\mathrm{geom}}^{\mathbf{T}}$ is the probability that the
drift bound~\eqref{eq:tsr-drift} is violated for some benign path $\gamma\in\mathcal{P}_S(L_0)$,
time $t$, and translator path $p$ of length at most the operating horizon;
\item for an MSR $\mathbf{M}$, $E_{\mathrm{geom}}^{\mathbf{M}}$ is the probability that the
distortion bound~\eqref{eq:msr-distortion} fails on some pair of induced manifold points
$z_\gamma(s),z_\gamma(t)$ coming from benign paths, or that curvature / chart regularity
assumptions used to define $\mathbf{M}$ are violated.
\end{itemize}
By construction, $E_{\mathrm{geom}}^{\mathbf{R}}$ depends only on $S$ and the ideal maps
$(\varphi_i, T_{ij})$ or $(\varphi_i,\psi_i,g)$, not on observation noise.

\item $E_{\mathrm{link}}^{\mathbf{R}}$ is the \emph{linkage / observation error}:
the probability that the detection pipeline fails \emph{even when geometry is ideal}.
It captures the effect of observation noise and feature misspecification:
\begin{itemize}
\item for a TSR, this includes deviations between noisy inputs $\tilde{x}_i$ and ideal
renders $\pi_i(u)$ that cause $\varphi_i(\tilde{x}_i)$ to move far enough from
$\varphi_i(\pi_i(u))$ to flip $\Sigma_{\mathbf{T}}$ across a decision boundary;
\item for an MSR, it includes analogous noise-induced shifts of manifold points
$z\in\mathcal{M}$ and any failure of $\Sigma_{\mathbf{M}}$ to track $d_g$ as intended.
\end{itemize}

\item $\xi^{\mathbf{R}}$ is a \emph{calibration error} term quantifying how far the
calibrated mapping from $\Sigma_{\mathbf{R}}$ to probabilities departs from the true
conditional $P(Y=1\mid \Sigma_{\mathbf{R}})$ in the high-risk region.

\item $\zeta^{\mathbf{R}}$ is an \emph{abstention failure} term: the probability that the
system fails to abstain when it should (e.g., when distortion audits, feature-range checks,
or support conditions indicate that $\mathbf{R}$ is being used outside its validated regime).

\item $\delta_{\mathrm{indep}}^{\mathbf{R}} \in [0,1]$ is an \emph{independence slack} that
upper-bounds the gap between the joint failure probability of the four error sources
(geometry, linkage, calibration, abstention) and the product of their marginals. It
quantifies deviations from an idealized independence model.

\item $\tau_{\mathrm{vac}}\in(0,1)$ is a \emph{vacuity threshold}. When the combined error
budget (e.g., a sum or multiplicative form in later theorems) exceeds $\tau_{\mathrm{vac}}$,
we declare the resulting correctness bound non-informative and treat the detector as
operating outside its validated regime.
\end{enumerate}
\end{definition}

\paragraph{Three-layer view.}
Definitions~\ref{def:semantic-structure}--\ref{def:sameness-detector} separate three distinct
contributors to detection performance:
\begin{itemize}[leftmargin=*]
\item \emph{Semantic drift} $\sigma_S^{\mathbf{M}}(L)$ (Definition~\ref{def:semantic-drift})
is how far latent entities move in geometry under benign evolution. It is a property of
$(S,\mathbf{M})$ and the path family.

\item \emph{Representation error} $\varepsilon_{\mathrm{trans}}$ and $\varepsilon_{\mathrm{dist}}$
(Definitions~\ref{def:tsr}--\ref{def:msr}) measure how imperfectly translators and charts
realize the intended semantics on ideal data. These feed into $E_{\mathrm{geom}}^{\mathbf{R}}$.

\item \emph{Observation noise and linkage error} are captured by $E_{\mathrm{link}}^{\mathbf{R}}$:
even if geometry were perfect, noisy or shifted observations can still flip decisions.
\end{itemize}
This three-layer decomposition is the backbone of the later Error Budget Transfer Theorem:
it lets us move guarantees between $\SamTrans(S)$ and $\SamGeom(S)$ while keeping semantic
drift, representation quality, and observation noise conceptually distinct.

\section{Realization categories and $\varepsilon$-equivalence}
\label{sec:realization-categories}

Section~\ref{sec:sameness-structure} introduced semantic sameness structures $S$ and two ways of realizing
them in concrete systems: translation-based realizations (TSRs, Definition~\ref{def:tsr}) and
manifold-based realizations (MSRs, Definition~\ref{def:msr}). In this section we package these
realizations into two categories,
\[
\SamTrans(S)
\quad\text{and}\quad
\SamGeom(S),
\]
equipped with coarse ``error gauges'' on objects, and we state an approximate equivalence between
well-behaved subcategories. This makes precise the idea that TSRs and MSRs are two
$\varepsilon$-equivalent \emph{categories of realizations} of the same underlying sameness problem $S$.

Throughout this section, $\|\cdot\|$ denotes the natural inner-product norm in the relevant feature
or ambient space, and $\mathcal{I}$ is the finite index set of modalities from
Definition~\ref{def:sameness-structure}. We fix a sameness structure $S$ once and for all and
suppress it from the notation when convenient.

\subsection{Category of translation-based realizations}
\label{subsec:samtrans}

We first define the category whose objects are translation-based realizations of $S$.

\begin{definition}[Category $\SamTrans(S)$]
\label{def:samtrans-category}
Fix a semantic sameness structure $S$ as in Definition~\ref{def:sameness-structure}, with latent
space $I$, modality index set $\mathcal{I}$, observation spaces $\{X_i\}_{i\in\mathcal{I}}$, rendering
maps $\{\pi_i : I \to X_i\}_{i\in\mathcal{I}}$, sameness relation $\approx$, and benign path family
$\mathcal{P}_S(L)$.

\paragraph{Objects.}
An object of $\SamTrans(S)$ is a translation-based semantic realization (TSR)
\[
\mathbf{T}
=
\Bigl(
S,\,
\{V_i^{\mathbf{T}}\}_{i\in\mathcal{I}},\,
\{\varphi_i^{\mathbf{T}} : X_i \to V_i^{\mathbf{T}}\}_{i\in\mathcal{I}},\,
\{T_{ij}^{\mathbf{T}} : V_i^{\mathbf{T}} \to V_j^{\mathbf{T}}\}_{i,j\in\mathcal{I}},\,
G_{\mathbf{T}},\,
\varepsilon_{\mathrm{trans}}^{\mathbf{T}}(\cdot)
\Bigr)
\]
as in Definition~\ref{def:tsr}. For categorical purposes we implicitly identify each feature space
$V_i^{\mathbf{T}}$ with a fixed Euclidean space $\mathbb{R}^{d_i}$ up to isometry, so that different
TSRs share a common coordinate system per modality.

\paragraph{Morphisms.}
Given two TSRs $\mathbf{T}, \mathbf{T}' \in \SamTrans(S)$, a morphism
\[
R : \mathbf{T} \to \mathbf{T}'
\]
with tolerances $(\eta_{\mathrm{enc}}, \eta_{\mathrm{trans}}) \ge 0$ is a family of reparameterizations
\[
R = \{R_i : V_i^{\mathbf{T}} \to V_i^{\mathbf{T}'}\}_{i\in\mathcal{I}}
\]
such that for all in-regime latent entities $u$ in a fixed subset $\Omega^{\mathrm{ideal}} \subseteq I$
(e.g., the union of points visited by benign paths in $\mathcal{P}_S(L_\star)$ for some operational
horizon $L_\star$) the following hold:
\begin{enumerate}[]
\item \textbf{Encoder compatibility.} The $R_i$ approximately intertwine the encoders:
\begin{equation}
\label{eq:tsr-enc-compat}
\bigl\|
R_i\bigl(\varphi_i^{\mathbf{T}}(\pi_i(u))\bigr)
-
\varphi_i^{\mathbf{T}'}(\pi_i(u))
\bigr\|
\;\le\;
\eta_{\mathrm{enc}}
\quad
\text{for all } i \in \mathcal{I}.
\end{equation}

\item \textbf{Translator compatibility.} The $R_i$ approximately commute with translators:
\begin{equation}
\label{eq:tsr-trans-compat}
\bigl\|
R_j\bigl(T_{ij}^{\mathbf{T}}(\varphi_i^{\mathbf{T}}(\pi_i(u)))\bigr)
-
T_{ij}^{\mathbf{T}'}\bigl(R_i(\varphi_i^{\mathbf{T}}(\pi_i(u)))\bigr)
\bigr\|
\;\le\;
\eta_{\mathrm{trans}}
\quad
\text{for all } i,j \in \mathcal{I}.
\end{equation}
\end{enumerate}

\paragraph{Composition and identities.}
If $R : \mathbf{T} \to \mathbf{T}'$ and $R' : \mathbf{T}' \to \mathbf{T}''$ are morphisms, their
composition $R' \circ R : \mathbf{T} \to \mathbf{T}''$ is defined componentwise:
\[
(R' \circ R)_i := R'_i \circ R_i, \quad i \in \mathcal{I}.
\]
The identity morphism on $\mathbf{T}$ is the family of identity maps
\[
\mathrm{id}_{\mathbf{T}} := \{\mathrm{id}_{V_i^{\mathbf{T}}}\}_{i\in\mathcal{I}}.
\]
These operations make $\SamTrans(S)$ a category.
\end{definition}

Intuitively, objects of $\SamTrans(S)$ are full translator systems built on a fixed sameness structure
$S$; morphisms are approximate changes of feature coordinates that preserve the encoders and
translators up to small tolerances. The in-regime set $\Omega^{\mathrm{ideal}}$ encodes where these
comparisons are intended to be meaningful.

\subsection{Category of manifold-based realizations}
\label{subsec:samgeom}

We now define the category whose objects are manifold-based realizations of $S$.

\begin{definition}[Category $\SamGeom(S)$]
\label{def:samgeom-category}
Fix the same sameness structure $S$ as above.

\paragraph{Objects.}
An object of $\SamGeom(S)$ is a manifold-based semantic realization (MSR)
\[
\mathbf{M}
=
\Bigl(
S,\,
\{V_i^{\mathbf{M}}\}_{i\in\mathcal{I}},\,
\{\varphi_i^{\mathbf{M}} : X_i \to V_i^{\mathbf{M}}\}_{i\in\mathcal{I}},\,
(\mathcal{M}^{\mathbf{M}}, g^{\mathbf{M}}),\,
\{\psi_i^{\mathbf{M}} : V_i^{\mathbf{M}} \to \mathcal{M}^{\mathbf{M}}\}_{i\in\mathcal{I}},\,
\varepsilon_{\mathrm{dist}}^{\mathbf{M}}(\cdot)
\Bigr)
\]
as in Definition~\ref{def:msr}. For each $i \in \mathcal{I}$ the smooth map
$\psi_i^{\mathbf{M}} : V_i^{\mathbf{M}} \to \mathcal{M}^{\mathbf{M}}$ is a parameterization whose image
$U_i^{\mathbf{M}} := \psi_i^{\mathbf{M}}(V_i^{\mathbf{M}})$ is an open subset of the manifold, and
\[
\chi_i^{\mathbf{M}} := (\psi_i^{\mathbf{M}})^{-1} : U_i^{\mathbf{M}} \to V_i^{\mathbf{M}}
\]
is the associated chart map. The distortion profile $\varepsilon_{\mathrm{dist}}^{\mathbf{M}}(\cdot)$ controls
how the chart metrics compare to the Riemannian metric $g^{\mathbf{M}}$.

Given a latent entity $u \in I$ for which modality $i$ observes $u$ (i.e., $\pi_i(u)$ is defined),
we write
\[
v_i^{\mathbf{M}}(u)
:= \varphi_i^{\mathbf{M}}(\pi_i(u)) \in V_i^{\mathbf{M}},
\qquad
z^{\mathbf{M}}(u)
:= \psi_i^{\mathbf{M}}\bigl(v_i^{\mathbf{M}}(u)\bigr) \in \mathcal{M}^{\mathbf{M}}.
\]
By consistency of the MSR, $z^{\mathbf{M}}(u)$ does not depend on the choice of modality $i$ when
multiple modalities observe the same $u$.

\paragraph{Morphisms.}
Given two MSRs $\mathbf{M}, \mathbf{M}' \in \SamGeom(S)$, a morphism
\[
H : \mathbf{M} \to \mathbf{M}'
\]
with tolerances $(\eta_{\mathrm{enc}}, \eta_{\mathrm{chart}}) \ge 0$ is a pair
\[
H = \bigl(h, R\bigr),
\qquad
h : \mathcal{M}^{\mathbf{M}} \to \mathcal{M}^{\mathbf{M}'},
\quad
R = \{R_i : V_i^{\mathbf{M}} \to V_i^{\mathbf{M}'}\}_{i\in\mathcal{I}},
\]
where $h$ is smooth and $\{R_i\}$ are reparameterizations of the feature spaces, satisfying two
compatibility conditions on an in-regime latent subset $\Omega^{\mathrm{ideal}} \subseteq I$ and the
corresponding in-regime feature sets
\[
\Omega^{\mathrm{reg}}_i
:=
\varphi_i^{\mathbf{M}}\bigl(\pi_i(\Omega^{\mathrm{ideal}})\bigr)
\subseteq V_i^{\mathbf{M}}.
\]

\begin{enumerate}[]
\item \textbf{Encoder compatibility.} For all $u \in \Omega^{\mathrm{ideal}}$ and $i \in \mathcal{I}$,
\begin{equation}
\label{eq:msr-enc-compat}
\bigl\|
R_i\bigl(\varphi_i^{\mathbf{M}}(\pi_i(u))\bigr)
-
\varphi_i^{\mathbf{M}'}(\pi_i(u))
\bigr\|
\;\le\;
\eta_{\mathrm{enc}}.
\end{equation}

\item \textbf{$h$-induced reparameterization and chart compatibility.}
For each modality $i \in \mathcal{I}$, define the \emph{$h$-induced reparameterization}
\[
R_h^i : V_i^{\mathbf{M}} \to V_i^{\mathbf{M}'},
\qquad
R_h^i(v)
:=
\bigl(\chi_i^{\mathbf{M}'} \circ h \circ \psi_i^{\mathbf{M}}\bigr)(v).
\]
This map transports a feature vector $v \in V_i^{\mathbf{M}}$ onto the manifold
$\mathcal{M}^{\mathbf{M}}$, applies $h$, and then pulls back to $V_i^{\mathbf{M}'}$ using the chart
of $\mathbf{M}'$.

We require that $R$ and the $h$-induced maps agree up to $\eta_{\mathrm{chart}}$ on in-regime
features:
\begin{equation}
\label{eq:msr-chart-compat}
\sup_{v \in \Omega^{\mathrm{reg}}_i}
\bigl\|
R_i(v) - R_h^i(v)
\bigr\|
\;\le\;
\eta_{\mathrm{chart}}
\quad
\text{for all } i \in \mathcal{I}.
\end{equation}
\end{enumerate}

\paragraph{Composition and identities.}
If $H = (h, R) : \mathbf{M} \to \mathbf{M}'$ and
$H' = (h', R') : \mathbf{M}' \to \mathbf{M}''$ are morphisms, their composition is
\[
H' \circ H
:=
\bigl(h' \circ h,\; \{R'_i \circ R_i\}_{i\in\mathcal{I}}\bigr).
\]
The identity morphism on $\mathbf{M}$ is
\[
\mathrm{id}_{\mathbf{M}} := (\mathrm{id}_{\mathcal{M}^{\mathbf{M}}}, \{\mathrm{id}_{V_i^{\mathbf{M}}}\}_{i\in\mathcal{I}}).
\]
These operations make $\SamGeom(S)$ a category.
\end{definition}

Intuitively, morphisms in $\SamGeom(S)$ consist of a manifold map $h$ (acting in state space) and
a family of feature-space maps $R_i$; the condition~\eqref{eq:msr-chart-compat} says that $R_i$
must be close, on in-regime representations, to the reparameterization induced by $h$ via the charts.
This separates the definition of the induced map $R_h^i$ from the compatibility requirement and
makes the geometry of morphisms transparent.

\subsection{Error gauges on realization categories}
\label{subsec:error-gauges}

The categories $\SamTrans(S)$ and $\SamGeom(S)$ contain many realizations that differ in ways
irrelevant to detection. To compare them coarsely, we introduce object-level \emph{error gauges}
that act as pseudo-metrics on the object sets. These gauges focus on the components that enter
the detection theory: drift/distortion profiles and geometric behavior on in-regime entities.

\begin{definition}[Error gauges on $\SamTrans(S)$ and $\SamGeom(S)$]
\label{def:error-gauges}
Fix an in-regime latent subset $\Omega^{\mathrm{ideal}} \subseteq I$, for example the union of points
visited by benign paths in $\mathcal{P}_S(L_\star)$, and let $\sigma_S^{\mathbf{M}}(L)$ be the semantic
drift profile from Definition~\ref{def:semantic-drift} whenever an MSR $\mathbf{M}$ exists.

\paragraph{Gauge on $\SamTrans(S)$.}
For two TSRs $\mathbf{T}, \mathbf{T}' \in \SamTrans(S)$, define
\begin{align}
\Delta_{\SamTrans}(\mathbf{T}, \mathbf{T}')
&:=
\sup_{k \in \mathbb{N}}
\bigl|
\varepsilon_{\mathrm{trans}}^{\mathbf{T}}(k)
-
\varepsilon_{\mathrm{trans}}^{\mathbf{T}'}(k)
\bigr|
\notag\\
&\quad+
\sup_{i,j \in \mathcal{I}}
\sup_{u \in \Omega^{\mathrm{ideal}}}
\bigl\|
T_{ij}^{\mathbf{T}}\bigl(v_i^{\mathbf{T}}(u)\bigr)
-
T_{ij}^{\mathbf{T}'}\bigl(v_i^{\mathbf{T}'}(u)\bigr)
\bigr\|,
\label{eq:gauge-samtrans}
\end{align}
where
$v_i^{\mathbf{T}}(u) := \varphi_i^{\mathbf{T}}(\pi_i(u))$ and
$v_i^{\mathbf{T}'}(u) := \varphi_i^{\mathbf{T}'}(\pi_i(u))$.
We tacitly identify $V_i^{\mathbf{T}}$ and $V_i^{\mathbf{T}'}$ with a common $\mathbb{R}^{d_i}$ up to
isometry, so that the norm in~\eqref{eq:gauge-samtrans} is well-defined.

\paragraph{Gauge on $\SamGeom(S)$.}
For two MSRs $\mathbf{M}, \mathbf{M}' \in \SamGeom(S)$, define
\begin{align}
\Delta_{\SamGeom}(\mathbf{M}, \mathbf{M}')
&:=
\sup_{\ell > 0}
\bigl|
\varepsilon_{\mathrm{dist}}^{\mathbf{M}}(\ell)
-
\varepsilon_{\mathrm{dist}}^{\mathbf{M}'}(\ell)
\bigr|
\notag\\
&\quad+
\sup_{u, u' \in \Omega^{\mathrm{ideal}}}
\Bigl|
d_{g^{\mathbf{M}}}\bigl(z^{\mathbf{M}}(u), z^{\mathbf{M}}(u')\bigr)
-
d_{g^{\mathbf{M}'}}\bigl(z^{\mathbf{M}'}(u), z^{\mathbf{M}'}(u')\bigr)
\Bigr|,
\label{eq:gauge-samgeom}
\end{align}
where $z^{\mathbf{M}}(u)$ and $z^{\mathbf{M}'}(u)$ are the induced manifold states defined above.
\end{definition}

Both $\Delta_{\SamTrans}$ and $\Delta_{\SamGeom}$ are symmetric, non-negative, and vanish when
the corresponding drift/distortion profiles and geometric behaviors agree exactly, but they may not
separate all distinct objects (they are pseudo-metrics rather than metrics). They deliberately
ignore implementation details that do not affect geometric behavior on in-regime entities.

\subsection{$\varepsilon$-equivalence of realization categories}
\label{subsec:epsilon-equivalence}

We next formalize the sense in which $\SamTrans(S)$ and $\SamGeom(S)$ will be treated as
approximately equivalent. The key idea is that we do \emph{not} expect strict categorical
equivalence, but rather an equivalence up to small slack in the error gauges.

\begin{definition}[$\varepsilon$-equivalence of categories of realizations]
\label{def:epsilon-equivalence}
Let $(\mathcal{C}, \Delta_{\mathcal{C}})$ and $(\mathcal{D}, \Delta_{\mathcal{D}})$ be categories equipped
with object-level pseudo-metrics (error gauges) on $\mathrm{Ob}(\mathcal{C})$ and
$\mathrm{Ob}(\mathcal{D})$, respectively. Let $\mathcal{C}^0 \subseteq \mathcal{C}$ and
$\mathcal{D}^0 \subseteq \mathcal{D}$ be full subcategories.

We say that $\mathcal{C}^0$ and $\mathcal{D}^0$ are \emph{$\varepsilon$-equivalent} if there exist
functors
\[
F : \mathcal{C}^0 \to \mathcal{D}^0,
\qquad
G : \mathcal{D}^0 \to \mathcal{C}^0,
\]
and non-negative slack functions $C_F, C_G$ such that:
\begin{enumerate}[]
\item For every object $C \in \mathrm{Ob}(\mathcal{C}^0)$,
\[
\Delta_{\mathcal{C}}\bigl(G(F(C)), C\bigr)
\;\le\;
C_G(\varepsilon_C),
\]
where $\varepsilon_C$ collects the relevant distortion parameters of $C$ (e.g.,
$\varepsilon_{\mathrm{trans}}^{\mathbf{T}}$, $\varepsilon_{\mathrm{dist}}^{F(\mathbf{T})}$, chart slack, etc.).
\item For every object $D \in \mathrm{Ob}(\mathcal{D}^0)$,
\[
\Delta_{\mathcal{D}}\bigl(F(G(D)), D\bigr)
\;\le\;
C_F(\varepsilon_D),
\]
with $\varepsilon_D$ defined analogously.
\item (\emph{First-order property.}) The slack functions are \emph{first order} in the distortion
parameters: when all distortion and slack terms (e.g., $\varepsilon_{\mathrm{trans}}$,
$\varepsilon_{\mathrm{dist}}$, chart-overlap slack $\delta_{\mathrm{chart}}$, translator asymmetry
$\delta_{\mathrm{asymm}}$) tend to zero, $C_F$ and $C_G$ also tend to zero.
\end{enumerate}
\end{definition}

In our setting, $\mathcal{C}^0$ and $\mathcal{D}^0$ will be well-behaved subcategories
$\SamTrans^0(S) \subseteq \SamTrans(S)$ and $\SamGeom^0(S) \subseteq \SamGeom(S)$ on which
smooth chartability and bounded-distortion conditions hold. The functors $F$ and $G$ will be the
forward and reverse constructions $F_S$ and $G_S$ that move between translators and manifolds.

\subsection{Structural equivalence theorem and implications}
\label{subsec:structural-theorem}

We are now ready to state the main structural theorem of this section in informal form. It says
that, under a smooth chartability assumption, the TSR and MSR design spaces for a fixed sameness
structure $S$ are $\varepsilon$-equivalent in the sense of Definition~\ref{def:epsilon-equivalence}.

\begin{theorem}[Realization categories are $\varepsilon$-equivalent (informal)]
\label{thm:structural-equivalence}
Fix a semantic sameness structure $S$ and an in-regime latent subset
$\Omega^{\mathrm{ideal}} \subseteq I$. There exist:
\begin{itemize}
\item well-behaved subcategories
\[
\SamTrans^0(S) \subseteq \SamTrans(S),
\qquad
\SamGeom^0(S) \subseteq \SamGeom(S),
\]
consisting of TSRs and MSRs that satisfy smooth chartability and bounded-distortion conditions;

\item functors
\[
F_S : \SamTrans^0(S) \to \SamGeom^0(S),
\qquad
G_S : \SamGeom^0(S) \to \SamTrans^0(S),
\]
that:
\begin{enumerate}[]
\item preserve the underlying sameness structure $S$ and feature spaces
$\{V_i, \varphi_i\}_{i\in\mathcal{I}}$;
\item send translators to manifolds by quotienting and gluing charts (Section~4);
\item send manifolds to translators by extracting chart transitions (Section~5);
\end{enumerate}

\item slack functions $C_F, C_G$ which are first order in the distortion and chartability
parameters (translator drift $\varepsilon_{\mathrm{trans}}$, metric distortion
$\varepsilon_{\mathrm{dist}}$, chart-overlap slack $\delta_{\mathrm{chart}}$, and translator asymmetry
$\delta_{\mathrm{asymm}}$),
\end{itemize}
such that:
\begin{align*}
\Delta_{\SamTrans}\bigl(\mathbf{T},\; G_S(F_S(\mathbf{T}))\bigr)
&\;\le\;
C_G\bigl(\varepsilon_{\mathrm{trans}}^{\mathbf{T}}, \varepsilon_{\mathrm{dist}}^{F_S(\mathbf{T})},
\delta_{\mathrm{chart}}(\mathbf{T},F_S(\mathbf{T})), \delta_{\mathrm{asymm}}(\mathbf{T})\bigr)
\quad\text{for all } \mathbf{T} \in \SamTrans^0(S),\\[0.25em]
\Delta_{\SamGeom}\bigl(\mathbf{M},\; F_S(G_S(\mathbf{M}))\bigr)
&\;\le\;
C_F\bigl(\varepsilon_{\mathrm{dist}}^{\mathbf{M}}, \varepsilon_{\mathrm{trans}}^{G_S(\mathbf{M})},
\delta_{\mathrm{chart}}(G_S(\mathbf{M}),\mathbf{M}), \delta_{\mathrm{asymm}}(G_S(\mathbf{M}))\bigr)
\quad\text{for all } \mathbf{M} \in \SamGeom^0(S).
\end{align*}

In particular, as all distortion and slack terms tend to zero, the round-trip errors in the gauges
$\Delta_{\SamTrans}$ and $\Delta_{\SamGeom}$ tend to zero: $\SamTrans^0(S)$ and $\SamGeom^0(S)$
are $\varepsilon$-equivalent categories of realizations of $S$.
\end{theorem}

The detailed constructions of $F_S$ and $G_S$ and the proof of
Theorem~\ref{thm:structural-equivalence} are given in Sections~4--5. At a conceptual level:

\begin{itemize}
\item For a fixed sameness structure $S$, TSRs and MSRs are two different \emph{realizations} of the
same problem: $\SamTrans^0(S)$ and $\SamGeom^0(S)$ are two design spaces for the same semantic
sameness structure.

\item The functors $F_S$ and $G_S$ witness an $\varepsilon$-equivalence between these design spaces,
showing that---under smooth chartability---they are dual up to controlled slack. An engineer who
starts from only $S$ can choose either a translator-based or manifold-based architecture, with the
assurance that the theoretical guarantees of the resulting detector are comparable.

\item Because the equivalence is functorial and comes with explicit slack bounds, \emph{theorems and
tools transfer}: risk bounds or stability results proven in the manifold view can be pulled back to
translator systems via $G_S$, and translator-side diagnostics (e.g., spectral audits on format graphs)
can be pushed forward to furnish new manifold-side checks via $F_S$. This transfer principle underlies
the Error Budget Transfer Theorem in Section~6.
\end{itemize}

\section{Smooth chartability and the functor \texorpdfstring{$F_S$}{F\_S}}
\label{sec:chartability-F}

Fix a semantic sameness structure $S$ and the associated categories
$\SamTrans(S)$ and $\SamGeom(S)$ from Section~\ref{sec:categories}.
Throughout this section, we work with ideal latent entities
$u \in \Omega^{\mathrm{ideal}} \subseteq I$ and their encoder outputs
$v_i(u) := \varphi_i(\pi_i(u)) \in V_i$ as in Section~\ref{sec:sameness-structure}.
We write $\|\cdot\|$ for the natural inner-product norm in the relevant
feature or ambient space.

Our goal is to construct a functor
\[
F_S : \SamTrans^0(S) \longrightarrow \SamGeom^0(S)
\]
on a well-behaved subcategory $\SamTrans^0(S) \subseteq \SamTrans(S)$,
sending a translation-based realization (TSR) $\mathbf{T}$ to a
manifold-based realization (MSR) $F_S(\mathbf{T})$, and a morphism
$R : \mathbf{T} \to \mathbf{T}'$ to a morphism
$F_S(R) : F_S(\mathbf{T}) \to F_S(\mathbf{T}')$.
Crucially, $F_S$ must be \emph{constructed} from the TSR data, not defined
by fiat.

We proceed in three steps:
\begin{enumerate}
  \item We build a canonical \emph{translator metric quotient}
        $(\widetilde{\mathcal{M}}^{\mathbf{T}}, \tilde d_{\mathbf{T}})$
        from any TSR $\mathbf{T}$ (Section~\ref{subsec:translator-metric-quotient}).
  \item We define an \emph{intrinsic} notion of smooth chartability for TSRs
        that depends only on $\mathbf{T}$ itself and its latent coverage
        (Section~\ref{subsec:intrinsic-chartability}),
        and use it to carve out the subcategory $\SamTrans^0(S)$.
  \item We impose a single, explicit assumption of \emph{Riemannian
        realizability} for these metric spaces and use it to define $F_S$
        on objects and morphisms (Section~\ref{subsec:riemannian-realizability}).
\end{enumerate}
This isolates the genuinely geometric hypothesis---that the
translator-induced metric space is quasi-isometric to a smooth manifold---
from the intrinsic regularity properties of the TSR itself.

\subsection{The translator metric quotient}
\label{subsec:translator-metric-quotient}

Let $\mathbf{T} \in \SamTrans(S)$ be a TSR with components
\[
\mathbf{T}
=
\Bigl(
  S,\ \{V_i^{\mathbf{T}}\}_{i \in \mathcal{I}},\ 
  \{\varphi_i^{\mathbf{T}}\}_{i \in \mathcal{I}},\
  \{T_{ij}^{\mathbf{T}}\}_{(i,j) \in E_{\mathbf{T}}},\
  G_{\mathbf{T}},\ \varepsilon_{\mathrm{trans}}^{\mathbf{T}}(\cdot)
\Bigr),
\]
as in Definition~\ref{def:tsr}.  For each modality $i \in \mathcal{I}$,
recall the in-regime feature subset
\[
\Omega_i^{\mathrm{reg}}(\mathbf{T})
:= 
\bigl\{
  v_i(u) = \varphi_i^{\mathbf{T}}(\pi_i(u)) :
  u \in \Omega^{\mathrm{ideal}},\ \pi_i(u)\ \text{defined}
\bigr\}
\subseteq V_i^{\mathbf{T}},
\]
and define the disjoint union
\[
X^{\mathbf{T}}
:=
\bigsqcup_{i \in \mathcal{I}}
\Omega_i^{\mathrm{reg}}(\mathbf{T}),
\qquad
x = (i, v_i) \in X^{\mathbf{T}}.
\]

\paragraph{Local edge costs.}
We endow $X^{\mathbf{T}}$ with two families of elementary \emph{edges}:
\begin{itemize}
  \item \emph{Within-chart edges.} For any $i \in \mathcal{I}$ and any
        $v_i, v_i' \in \Omega_i^{\mathrm{reg}}(\mathbf{T})$ we introduce
        an edge $e : (i, v_i) \to (i, v_i')$ with cost
        \[
        c_{\mathrm{intra}}(e)
        :=
        \|v_i - v_i'\|.
        \]
  \item \emph{Translator edges.} For any $(i,j) \in E_{\mathbf{T}}$ and
        any $v_i \in \Omega_i^{\mathrm{reg}}(\mathbf{T})$ such that
        $T_{ij}^{\mathbf{T}}(v_i)$ lies in $\Omega_j^{\mathrm{reg}}(\mathbf{T})$,
        we introduce an edge
        \[
        e : (i, v_i) \longrightarrow (j, T_{ij}^{\mathbf{T}}(v_i))
        \]
        with cost
        \[
        c_{\mathrm{trans}}(e)
        :=
        \bigl\| T_{ij}^{\mathbf{T}}(v_i) \bigr\|.
        \]
        (Any consistent bounded positive edge-weighting based on the
        translator action would suffice; the specific choice is not
        essential for the subsequent theory.)
\end{itemize}

\paragraph{Path pseudo-metric on $X^{\mathbf{T}}$.}
A \emph{finite edge-path} in $X^{\mathbf{T}}$ is a sequence
$p = (e_1, \dots, e_m)$ of composable edges.
We define its cost
\[
\mathrm{len}(p)
:=
\sum_{r=1}^m c(e_r),
\]
where $c(e)$ is $c_{\mathrm{intra}}(e)$ or $c_{\mathrm{trans}}(e)$
depending on the edge type.
For any two points $x,y \in X^{\mathbf{T}}$, we define the
\emph{path pseudo-metric}
\[
\tilde d_{\mathbf{T}}^{\mathrm{raw}}(x,y)
:=
\inf\bigl\{ \mathrm{len}(p) : p\ \text{is a finite edge-path from $x$ to $y$} \bigr\}.
\]
Standard arguments show that $\tilde d_{\mathbf{T}}^{\mathrm{raw}}$ is a
pseudo-metric: it is symmetric, nonnegative, and satisfies the triangle
inequality, but may assign zero distance to distinct points.

\begin{definition}[Translator metric quotient]
\label{def:translator-metric-quotient}
Define an equivalence relation $\sim_0$ on $X^{\mathbf{T}}$ by
\[
x \sim_0 y
\quad\Longleftrightarrow\quad
\tilde d_{\mathbf{T}}^{\mathrm{raw}}(x,y) = 0.
\]
The \emph{translator metric quotient} of $\mathbf{T}$ is the metric space
\[
\bigl(
  \widetilde{\mathcal{M}}^{\mathbf{T}},
  \tilde d_{\mathbf{T}}
\bigr)
\]
obtained as the quotient $X^{\mathbf{T}} / \sim_0$ with metric
\[
\tilde d_{\mathbf{T}}([x], [y])
:=
\tilde d_{\mathbf{T}}^{\mathrm{raw}}(x,y),
\]
which is well-defined and satisfies the metric axioms.

For each modality $i \in \mathcal{I}$ there is a canonical
\emph{quotient parameterization}
\[
\tilde\psi_i^{\mathbf{T}} : \Omega_i^{\mathrm{reg}}(\mathbf{T}) \to
\widetilde{\mathcal{M}}^{\mathbf{T}},
\qquad
\tilde\psi_i^{\mathbf{T}}(v_i) := [(i, v_i)].
\]
\end{definition}

Intuitively, $\widetilde{\mathcal{M}}^{\mathbf{T}}$ is the ``patchwork
quilt'' you obtain by gluing together the in-regime feature charts
along the translator edges, with $\tilde d_{\mathbf{T}}$ measuring the
shortest-path cost induced by chart distances and translator hops.
The construction is canonical (up to isometry) and depends only on the
TSR $\mathbf{T}$.

\subsection{Intrinsic smooth chartability}
\label{subsec:intrinsic-chartability}

The previous construction associates a metric space
$(\widetilde{\mathcal{M}}^{\mathbf{T}}, \tilde d_{\mathbf{T}})$ to any
TSR $\mathbf{T}$.  However, for an arbitrary $\mathbf{T}$ this space
may be pathological: it need not resemble a manifold at any scale.
We now isolate a set of purely \emph{intrinsic} conditions on
$\mathbf{T}$ under which the translator metric quotient behaves like
a well-controlled, finite-dimensional geometric object.

\begin{definition}[Intrinsic smooth chartability]
\label{def:intrinsic-chartability}
A TSR $\mathbf{T}$ is \emph{intrinsically smoothly chartable} if the
following conditions hold on an in-regime subset $\Omega^{\mathrm{ideal}}$
and the associated feature sets $\Omega_i^{\mathrm{reg}}(\mathbf{T})$:
\begin{enumerate}[]
  \item \textbf{Translator regularity.} For each $(i,j) \in E_{\mathbf{T}}$,
        the translator $T_{ij}^{\mathbf{T}} : V_i^{\mathbf{T}} \to
        V_j^{\mathbf{T}}$ is $C^2$ on an open neighbourhood of
        $\Omega_i^{\mathrm{reg}}(\mathbf{T})$, with uniformly bounded
        Jacobian and inverse Jacobian:
        \[
        \sup_{v_i \in \Omega_i^{\mathrm{reg}}(\mathbf{T})}
        \bigl\|
          D T_{ij}^{\mathbf{T}}(v_i)
        \bigr\|
        \le C_{\mathrm{Jac}},
        \qquad
        \sup_{v_i \in \Omega_i^{\mathrm{reg}}(\mathbf{T})}
        \bigl\|
          D T_{ij}^{\mathbf{T}}(v_i)^{-1}
        \bigr\|
        \le C_{\mathrm{Jac}}.
        \]
  \item \textbf{Round-trip asymmetry bound.} There exists
        $\delta_{\mathrm{asymm}}^{\mathbf{T}} < \infty$ such that for every
        $(i,j) \in E_{\mathbf{T}}$ and every
        $v_i \in \Omega_i^{\mathrm{reg}}(\mathbf{T})$ with
        $T_{ij}^{\mathbf{T}}(v_i) \in \Omega_j^{\mathrm{reg}}(\mathbf{T})$
        and $T_{ji}^{\mathbf{T}}(T_{ij}^{\mathbf{T}}(v_i))
        \in \Omega_i^{\mathrm{reg}}(\mathbf{T})$ we have
        \[
        \bigl\|
          T_{ji}^{\mathbf{T}}(T_{ij}^{\mathbf{T}}(v_i)) - v_i
        \bigr\|
        \le \delta_{\mathrm{asymm}}^{\mathbf{T}}.
        \]
  \item \textbf{Cocycle bound.} There exists
        $\delta_{\mathrm{cocycle}}^{\mathbf{T}} < \infty$ such that for any
        triple $(i,j,k)$ with $(i,j), (j,k), (i,k) \in E_{\mathbf{T}}$ and
        any $v_i \in \Omega_i^{\mathrm{reg}}(\mathbf{T})$ lying in the domain
        of all three compositions,
        \[
        \bigl\|
          T_{jk}^{\mathbf{T}}(T_{ij}^{\mathbf{T}}(v_i)) -
          T_{ik}^{\mathbf{T}}(v_i)
        \bigr\|
        \le \delta_{\mathrm{cocycle}}^{\mathbf{T}}.
        \]
        This ensures that multi-hop translator paths are consistent, up to a
        controlled slack, with direct translations.
  \item \textbf{Latent coverage and drift control.}
        There is an operational horizon $L_\star > 0$ such that:
        \begin{itemize}
          \item every benign latent path $\gamma \in \mathcal{P}_S(L_\star)$
                projects to in-regime observations $\pi_i(\gamma(t))$ and
                hence to in-regime features
                $v_i(\gamma(t)) = \varphi_i^{\mathbf{T}}(\pi_i(\gamma(t)))$;
          \item the drift profile $\varepsilon_{\mathrm{trans}}^{\mathbf{T}}(k)$
                is finite for all $k$ and satisfies the per-time-step bound
                of Definition~\ref{def:tsr} on these paths.
        \end{itemize}
\end{enumerate}
\end{definition}

\begin{definition}[Well-behaved subcategory $\SamTrans^0(S)$]
\label{def:samtrans0}
We define $\SamTrans^0(S)$ as the full subcategory of
$\SamTrans(S)$ whose objects are intrinsically smoothly chartable
TSRs in the sense of Definition~\ref{def:intrinsic-chartability}, and whose
morphisms are exactly the morphisms of $\SamTrans(S)$ between
such objects.
\end{definition}

By construction, membership in $\SamTrans^0(S)$ is a property of the TSR
$\mathbf{T}$ and its translators alone.  No manifold is mentioned in
Definition~\ref{def:intrinsic-chartability}: it is a falsifiable hypothesis
about the regularity, consistency, and coverage of the translators.

\subsection{Riemannian realizability of the translator metric}
\label{subsec:riemannian-realizability}

The translator metric quotient
$(\widetilde{\mathcal{M}}^{\mathbf{T}}, \tilde d_{\mathbf{T}})$ from
Definition~\ref{def:translator-metric-quotient} is a canonical metric space
built from $\mathbf{T}$ alone.  Intrinsic chartability ensures that this
space behaves like a finite-dimensional, locally well-behaved metric
space: translators are smooth with bounded Jacobians, roundtrip and
cocycle inconsistencies are small, and latent data populate a connected,
bounded region.

To bridge from metric geometry to Riemannian geometry we isolate the
following explicit assumption, which is the only place where we posit
the existence of a smooth manifold ``shadow.''

\begin{assumption}[Riemannian realizability of translator metric]
\label{assump:riemannian-realizability}
For every $\mathbf{T} \in \SamTrans^0(S)$, the translator metric
quotient $(\widetilde{\mathcal{M}}^{\mathbf{T}}, \tilde d_{\mathbf{T}})$
admits a Riemannian realization: there exists a smooth, finite-dimensional
Riemannian manifold $(\mathcal{M}^{\mathbf{T}}, g^{\mathbf{T}})$ and
a quasi-isometry
\[
q^{\mathbf{T}} :
\bigl(\widetilde{\mathcal{M}}^{\mathbf{T}}, \tilde d_{\mathbf{T}}\bigr)
\longrightarrow
\bigl(\mathcal{M}^{\mathbf{T}}, d_{g^{\mathbf{T}}}\bigr)
\]
with distortion bounded by a function
$\varepsilon_{\mathrm{realize}}^{\mathbf{T}}(\cdot)$ that depends
continuously and monotonically on the intrinsic slack parameters
$\varepsilon_{\mathrm{trans}}^{\mathbf{T}}(\cdot)$,
$\delta_{\mathrm{asymm}}^{\mathbf{T}}$ and
$\delta_{\mathrm{cocycle}}^{\mathbf{T}}$.

Moreover, we assume that $q^{\mathbf{T}}$ is coarse-biLipschitz on the
in-regime region populated by $\Omega^{\mathrm{ideal}}$, so that
for all $x,y$ in this region
\[
\Bigl|
  \frac{d_{g^{\mathbf{T}}}\bigl(q^{\mathbf{T}}(x), q^{\mathbf{T}}(y)\bigr)}
       {\tilde d_{\mathbf{T}}(x,y)}
  - 1
\Bigr|
\le
\varepsilon_{\mathrm{realize}}^{\mathbf{T}}
\bigl(\tilde d_{\mathbf{T}}(x,y)\bigr).
\]
\end{assumption}

We do \emph{not} attempt to prove
Assumption~\ref{assump:riemannian-realizability} in this work; it is our
single, explicit geometric hypothesis.  In practice, it is an empirical
question whether a given TSR $\mathbf{T}$ satisfies this assumption,
and Section~\ref{sec:audits} discusses audit procedures for probing
Riemannian realizability.

\subsection{Constructing \texorpdfstring{$F_S$}{F\_S} on objects}
\label{subsec:F-on-objects}

With the metric quotient and realizability assumption in hand, we can
now define $F_S$ on objects in a non-circular way.

\begin{definition}[Functor $F_S$ on objects]
\label{def:F-on-objects}
Let $\mathbf{T} \in \SamTrans^0(S)$ be an intrinsically smoothly chartable
TSR.  Consider its translator metric quotient
$(\widetilde{\mathcal{M}}^{\mathbf{T}}, \tilde d_{\mathbf{T}})$ and a
Riemannian realization
$(\mathcal{M}^{\mathbf{T}}, g^{\mathbf{T}})$ with quasi-isometry
$q^{\mathbf{T}}$ given by Assumption~\ref{assump:riemannian-realizability}.
For each modality $i \in \mathcal{I}$, define the parameterization
\[
\psi_i^{\mathbf{T}}
: V_i^{\mathbf{T}} \supseteq \Omega_i^{\mathrm{reg}}(\mathbf{T})
\longrightarrow
\mathcal{M}^{\mathbf{T}},
\qquad
\psi_i^{\mathbf{T}}(v_i)
:=
q^{\mathbf{T}}\bigl(\tilde\psi_i^{\mathbf{T}}(v_i)\bigr),
\]
where $\tilde\psi_i^{\mathbf{T}}$ is the quotient map from
Definition~\ref{def:translator-metric-quotient}.

We define
\[
F_S(\mathbf{T})
:=
\mathbf{M}^{\mathbf{T}}
:=
\Bigl(
  S,\
  \{V_i^{\mathbf{T}}\}_{i\in\mathcal{I}},\
  \{\varphi_i^{\mathbf{T}}\}_{i\in\mathcal{I}},\
  (\mathcal{M}^{\mathbf{T}}, g^{\mathbf{T}}),\
  \{\psi_i^{\mathbf{T}}\}_{i\in\mathcal{I}},\
  \varepsilon_{\mathrm{dist}}^{\mathbf{T}}(\cdot)
\Bigr),
\]
where the distortion profile $\varepsilon_{\mathrm{dist}}^{\mathbf{T}}$ is defined
by pulling back the quasi-isometry distortion:
\[
\Bigl|
  \frac{\delta_i^{\mathbf{T}}(z_a,z_b)}
       {d_{g^{\mathbf{T}}}(z_a,z_b)}
  - 1
\Bigr|
\le
\varepsilon_{\mathrm{dist}}^{\mathbf{T}}\bigl(d_{g^{\mathbf{T}}}(z_a,z_b)\bigr),
\qquad
z_a, z_b \in \psi_i^{\mathbf{T}}(\Omega_i^{\mathrm{reg}}(\mathbf{T})),
\]
with chart distances
$\delta_i^{\mathbf{T}}(z_a,z_b) :=
\|\chi_i^{\mathbf{T}}(z_a) - \chi_i^{\mathbf{T}}(z_b)\|$
and chart maps $\chi_i^{\mathbf{T}} := (\psi_i^{\mathbf{T}})^{-1}$ on their
images.

By construction, $\mathbf{M}^{\mathbf{T}}$ is an MSR of $S$ in the sense
of Definition~\ref{def:msr}, and we have $F_S(\mathbf{T}) \in \SamGeom^0(S)$.
\end{definition}

Importantly, $F_S(\mathbf{T})$ is now determined (up to quasi-isometry) by
$\mathbf{T}$ itself: the only freedom lies in the choice of Riemannian
realization $(\mathcal{M}^{\mathbf{T}}, g^{\mathbf{T}})$ of the canonical
metric quotient.  Different choices of $q^{\mathbf{T}}$ give rise to
MSRs that are close in the error gauge $\Delta_{\SamGeom}$ of
Definition~\ref{def:error-gauges}, and Theorem~\ref{thm:structural-equivalence}
will show that this slack is first-order in the intrinsic errors of
$\mathbf{T}$.

\subsection{Constructing \texorpdfstring{$F_S$}{F\_S} on morphisms}
\label{subsec:F-on-morphisms}

We next define $F_S$ on morphisms.  Let
$R : \mathbf{T} \to \mathbf{T}'$ be a morphism in
$\SamTrans^0(S)$ as in Definition~\ref{def:samtrans-category},
with reparameterizations
$R_i : V_i^{\mathbf{T}} \to V_i^{\mathbf{T}'}$ satisfying the
encoder- and translator-compatibility bounds on ideal data.

The TSR morphism $R$ induces a map between the translator metric
quotients,
\[
\tilde h^R :
\widetilde{\mathcal{M}}^{\mathbf{T}}
\longrightarrow
\widetilde{\mathcal{M}}^{\mathbf{T}'},
\]
defined on representatives by
\[
\tilde h^R\bigl(\tilde\psi_i^{\mathbf{T}}(v_i)\bigr)
:=
\tilde\psi_i^{\mathbf{T}'}\bigl(R_i(v_i)\bigr),
\]
and extended by continuity.  The compatibility conditions in
Definition~\ref{def:samtrans-category} ensure that $\tilde h^R$ is
well-defined up to the metric identification $\sim_0$ and is
Lipschitz (with constants depending on the morphism error
tolerances $(\eta_{\mathrm{enc}}, \eta_{\mathrm{trans}})$).

Using the quasi-isometries $q^{\mathbf{T}}$ and $q^{\mathbf{T}'}$ from
Assumption~\ref{assump:riemannian-realizability}, we then define the
manifold-level map
\[
h^R :
\mathcal{M}^{\mathbf{T}}
\longrightarrow
\mathcal{M}^{\mathbf{T}'},
\qquad
h^R
:=
q^{\mathbf{T}'}
\circ \tilde h^R
\circ (q^{\mathbf{T}})^{-1},
\]
understanding $(q^{\mathbf{T}})^{-1}$ as a coarse inverse on the
in-regime region.

\begin{definition}[Functor $F_S$ on morphisms]
\label{def:F-on-morphisms}
For a morphism $R : \mathbf{T} \to \mathbf{T}'$ in
$\SamTrans^0(S)$, we define
\[
F_S(R)
:=
H^R
:=
\bigl(h^R,\ R\bigr)
:
\mathbf{M}^{\mathbf{T}}
\longrightarrow
\mathbf{M}^{\mathbf{T}'},
\]
where $\mathbf{M}^{\mathbf{T}} = F_S(\mathbf{T})$ and
$\mathbf{M}^{\mathbf{T}'} = F_S(\mathbf{T}')$, and $h^R$ is the
manifold map defined above.

The encoder-compatibility condition for $H^R$,
\[
\bigl\|
  R_i\bigl(\varphi_i^{\mathbf{M}^{\mathbf{T}}}(\pi_i(u))\bigr)
  -
  \varphi_i^{\mathbf{M}^{\mathbf{T}'}}(\pi_i(u))
\bigr\|
\le \eta_{\mathrm{enc}}',
\]
is inherited from $R$, while the chart-compatibility condition,
\[
\bigl\|
  R_i\bigl(\varphi_i^{\mathbf{M}^{\mathbf{T}}}(\pi_i(u))\bigr)
  -
  \chi_i^{\mathbf{M}^{\mathbf{T}'}}\bigl(h^R(z^{\mathbf{T}}(u))\bigr)
\bigr\|
\le \eta_{\mathrm{chart}}',
\]
with $z^{\mathbf{T}}(u)$ the latent state in $\mathcal{M}^{\mathbf{T}}$,
follows from the construction of $h^R$ and the quasi-isometry bounds.
Thus $H^R$ is a morphism in $\SamGeom(S)$ in the sense of
Definition~\ref{def:samgeom-category}.
\end{definition}

\begin{proposition}[Functoriality and stability of $F_S$]
\label{prop:F-functorial}
The assignment $F_S$ of Definitions~\ref{def:F-on-objects}
and~\ref{def:F-on-morphisms} defines a functor
\[
F_S : \SamTrans^0(S) \longrightarrow \SamGeom^0(S)
\]
satisfying:
\begin{enumerate}[]
  \item For every object $\mathbf{T} \in \SamTrans^0(S)$,
        $F_S(\mathrm{id}_{\mathbf{T}}) = \mathrm{id}_{F_S(\mathbf{T})}$.
  \item For any composable morphisms
        $R : \mathbf{T} \to \mathbf{T}'$ and
        $R' : \mathbf{T}' \to \mathbf{T}''$ in $\SamTrans^0(S)$,
        \[
        F_S(R' \circ R)
        =
        F_S(R') \circ F_S(R).
        \]
  \item There exists a universal, increasing function $C_F$
        such that for any $\mathbf{T}, \mathbf{T}' \in \SamTrans^0(S)$
        the error gauge satisfies
        \[
        \Delta_{\SamGeom}\bigl(F_S(\mathbf{T}), F_S(\mathbf{T}')\bigr)
        \le
        C_F\Bigl(
          \Delta_{\SamTrans}(\mathbf{T}, \mathbf{T}'),\
          \delta_{\mathrm{chart}}(\mathbf{T}),\
          \delta_{\mathrm{chart}}(\mathbf{T}')
        \Bigr),
        \]
        where $\delta_{\mathrm{chart}}(\mathbf{T})$ is the chart-overlap
        slack of $\mathbf{T}$ from Definition~\ref{def:chart-slack}.
        Moreover, $C_F$ is first-order in its arguments: it vanishes as
        $\Delta_{\SamTrans} \to 0$ and $\delta_{\mathrm{chart}} \to 0$.
\end{enumerate}
\end{proposition}

\begin{proof}[Proof sketch]
Functoriality of $F_S$ on objects and morphisms follows from the
definitions of the translator metric quotient and the functoriality of
the quasi-isometry construction in Assumption~\ref{assump:riemannian-realizability}.
The stability bound in (iii) leverages the fact that the distortion
profile $\varepsilon_{\mathrm{dist}}^{\mathbf{T}}$ is controlled by the
intrinsic errors of $\mathbf{T}$ via the quasi-isometry, and that the
error gauge $\Delta_{\SamGeom}$ is built from these distortion profiles
and chart disagreements.  Full details are deferred to
Section~6, where we prove the stronger Error Budget Transfer Theorem.
\end{proof}

\subsection{Discussion: what $F_S$ really does}

The functor $F_S$ no longer hides a manifold inside the definition of
smooth chartability.  Instead, it proceeds in two conceptually clean
stages:

\begin{itemize}
  \item \textbf{From flat maps to a quilt.}  For any intrinsically
        well-behaved TSR $\mathbf{T}$, we canonically construct the
        translator metric quotient
        $(\widetilde{\mathcal{M}}^{\mathbf{T}}, \tilde d_{\mathbf{T}})$
        that encodes how features in different modalities can be
        transported and compared via translators.  This is the ``patchwork
        quilt'' of all flat maps stitched together along their edges,
        with a well-defined metric.
  \item \textbf{From the quilt to the globe.}  The single geometric
        hypothesis of Riemannian realizability asserts that this quilt
        is quasi-isometric to some smooth manifold
        $(\mathcal{M}^{\mathbf{T}}, g^{\mathbf{T}})$.  The functor $F_S$
        then declares this manifold, together with the induced
        parameterizations $\psi_i^{\mathbf{T}}$, to be the MSR
        $F_S(\mathbf{T})$.
\end{itemize}

In this way, $F_S$ is a genuine construction rather than a renaming:
it builds a metric geometry from the translators, and only then---under
an explicit, isolated assumption---passes to a smooth Riemannian
geometry.  The error gauge stability in
Proposition~\ref{prop:F-functorial} will feed directly into the Error
Budget Transfer Theorem in Section~6, where we show that detection
guarantees can be transported across this functor with only
first-order slack.
\section{From manifolds back to translators: the functor \texorpdfstring{$G_S$}{G\_S}}
\label{sec:G-construction}

Section~\ref{sec:chartability-F} constructed a functor
$F_S : \SamTrans^0(S) \to \SamGeom^0(S)$
by starting from a smoothly chartable TSR~$\mathbf{T}$, forming its translator
metric quotient $(\widetilde{\mathcal{M}}^{\mathbf{T}}, \tilde d_{\mathbf{T}})$, and---under a
Riemannian realizability assumption---obtaining a manifold-based realization
$F_S(\mathbf{T})$ whose charts approximate the original translators.
In this section we build the reverse functor
\[
G_S : \SamGeom^0(S) \to \SamTrans^0(S),
\]
which takes a well-parameterized manifold-based realization~$\mathbf{M}$ and
extracts a canonical translator-based realization $G_S(\mathbf{M})$. At the
level of the Globe–Flat Maps analogy, $F_S$ turns a consistent atlas of flat
maps into a globe; $G_S$ takes a globe with well-behaved charts and recovers
a translator graph encoding the chart transitions.

Throughout this section we fix a semantic sameness structure $S$ and work
within the realization categories $\SamGeom(S)$ and $\SamTrans(S)$
from Section~\ref{sec:categories}. Norms $\|\cdot\|$ denote the natural inner-product
norms on the relevant feature spaces, as in Section~\ref{sec:categories}.

\subsection{Geometrically well-parameterized MSRs}
\label{subsec:geom-well-param}

We first specify which manifold-based realizations admit a stable translation
view. Intuitively, an MSR is \emph{geometrically well-parameterized} when its
chart maps behave like well-conditioned coordinate systems on the region
traced out by benign latent paths, with controlled distortion and overlap.

\begin{definition}[Geometrically well-parameterized MSR]
\label{def:geom-well-param}
Let $\mathbf{M}$ be an MSR of $S$ in the sense of
Definition~\ref{def:msr}:
\[
\mathbf{M}
=
\Bigl(
S,\ \{V_i^{\mathbf{M}}\}_{i\in\mathcal{I}},\ \{\varphi_i^{\mathbf{M}}\}_{i\in\mathcal{I}},
(\mathcal{M}^{\mathbf{M}}, g^{\mathbf{M}}),\ \{\psi_i^{\mathbf{M}}\}_{i\in\mathcal{I}},
\varepsilon_{\mathrm{dist}}^{\mathbf{M}}(\cdot)
\Bigr).
\]
Write $U_i^{\mathbf{M}} := \psi_i^{\mathbf{M}}(V_i^{\mathbf{M}})$ and
$\chi_i^{\mathbf{M}} := (\psi_i^{\mathbf{M}})^{-1} : U_i^{\mathbf{M}} \to V_i^{\mathbf{M}}$,
and define $z^{\mathbf{M}}(u)$ as in Definition~\ref{def:semantic-drift}.

We say that $\mathbf{M}$ is \emph{geometrically well-parameterized} if the
following conditions hold.

\begin{enumerate}[label=(\roman*),leftmargin=*]

\item \textbf{Shared feature model.}
For each modality $i\in\mathcal{I}$ the feature spaces and encoders agree
(up to isometry) with those used in the TSR view:
there exist inner-product-preserving isomorphisms
$Q_i : V_i^{\mathbf{M}} \to V_i$ such that, after identifying
$V_i^{\mathbf{M}}$ with $V_i$ via $Q_i$, we may write
$\varphi_i^{\mathbf{M}} = \varphi_i$ and use the common notation
$V_i$, $\varphi_i$.

\item \textbf{Chart regularity and bounded distortion.}
Each parameterization $\psi_i : V_i \to \mathcal{M}^{\mathbf{M}}$ is $C^2$
on an in-regime region $V_i^{\mathrm{reg}} \subseteq V_i$ that contains the
image of the benign observation set under $\varphi_i \circ \pi_i$.
The corresponding chart maps
$\chi_i : U_i \to V_i$ with $U_i := \psi_i(V_i^{\mathrm{reg}})$ satisfy the
distortion profile of Definition~\ref{def:msr}: for all
$z_a,z_b \in U_i$ with $d_{g^{\mathbf{M}}}(z_a,z_b) > 0$,
\begin{equation}
\label{eq:chart-distortion}
\left|
\frac{\delta_i(z_a,z_b)}{d_{g^{\mathbf{M}}}(z_a,z_b)} - 1
\right|
\le
\varepsilon_{\mathrm{dist}}^{\mathbf{M}}\bigl(d_{g^{\mathbf{M}}}(z_a,z_b)\bigr),
\qquad
\delta_i(z_a,z_b) := \bigl\|\chi_i(z_a) - \chi_i(z_b)\bigr\|_{V_i}.
\end{equation}

\item \textbf{Chart overlap control.}
For any $i,j$ with $U_i \cap U_j \neq \varnothing$ and any
$z \in U_i \cap U_j$,
the local chart transition
$\chi_j \circ \psi_i : V_i^{\mathrm{reg}} \to V_j$
is well-defined on a neighborhood of $\chi_i(z)$ and has uniformly bounded
Jacobian and inverse Jacobian there. Equivalently, the family
$\{\chi_i\}$ forms a uniformly well-conditioned atlas on the in-regime
region of $\mathcal{M}^{\mathbf{M}}$.

\item \textbf{Benign path coverage.}
There exists an operational horizon $L_\star$ such that for every benign
latent path $\gamma \in \mathcal{P}_S(L_\star)$, the induced manifold path
$z^{\mathbf{M}}(\cdot)$ remains in $\bigcup_i U_i$ and crosses only finitely
many chart overlaps. The semantic drift profile
$\sigma_S^{\mathbf{M}}(L_\star)$ from Definition~\ref{def:semantic-drift}
is finite.

\end{enumerate}
\end{definition}

Condition~(i) ensures that the manifold view and the translator view share
the same feature model; (ii)–(iii) guarantee that chart coordinates are
well-behaved and that chart transitions are controlled; and (iv) ensures
that the benign path family never leaves the region where charts are valid.
Taken together, these conditions are the geometric mirror of the smooth
chartability conditions imposed on TSRs in
Section~\ref{sec:chartability-F}.

\begin{definition}[Well-behaved geometric subcategory]
\label{def:samgeom0}
Let $\SamGeom^0(S)$ denote the full subcategory of $\SamGeom(S)$ whose
objects are geometrically well-parameterized MSRs in the sense of
Definition~\ref{def:geom-well-param}, and whose morphisms are those of
$\SamGeom(S)$ (Definition~\ref{def:samgeom-category}) between such objects.
\end{definition}

\subsection{Constructing \texorpdfstring{$G_S$}{G\_S} on objects}
\label{subsec:G-on-objects}

We now define $G_S$ on objects by ``unfolding'' a well-parameterized manifold
back into a TSR whose translators are given by chart transitions.

\begin{definition}[Functor $G_S$ on objects]
\label{def:G-on-objects}
Let $\mathbf{M} \in \SamGeom^0(S)$ with data as in
Definition~\ref{def:geom-well-param}. We define the TSR
$G_S(\mathbf{M}) =: \mathbf{T}^{\mathbf{M}}$ by
\[
\mathbf{T}^{\mathbf{M}}
=
\Bigl(
S,\ \{V_i^{\mathbf{M}}\}_{i\in\mathcal{I}},\ \{\varphi_i^{\mathbf{M}}\}_{i\in\mathcal{I}},
\{T_{ij}^{\mathbf{M}}\}_{(i,j)\in E^{\mathbf{M}}},\ G_{\mathbf{T}^{\mathbf{M}}},
\varepsilon_{\mathrm{trans}}^{\mathbf{M}}(\cdot)
\Bigr),
\]
where:
\begin{enumerate}[label=(\roman*),leftmargin=*]

\item \textbf{Feature spaces and encoders.}
We reuse the feature spaces and encoders of $\mathbf{M}$:
$V_i^{\mathbf{M}}$ and $\varphi_i^{\mathbf{M}}$ become the feature spaces and
encoders of $\mathbf{T}^{\mathbf{M}}$.

\item \textbf{Translator graph.}
The vertex set of $G_{\mathbf{T}^{\mathbf{M}}}$ is $\mathcal{I}$, and we include
a directed edge $(i,j)$ whenever $U_i \cap U_j$ is non-empty on the in-regime
region traced by benign paths. Thus translators exist exactly where chart
overlaps do.

\item \textbf{Chart-transition translators.}
For each edge $(i,j) \in E^{\mathbf{M}}$ we define the translator
$T_{ij}^{\mathbf{M}} : V_i^{\mathbf{M}} \to V_j^{\mathbf{M}}$ on the in-regime
domain $V_i^{\mathrm{reg}}$ by
\begin{equation}
\label{eq:chart-translator}
T_{ij}^{\mathbf{M}}(v_i)
:=
\chi_j^{\mathbf{M}}\bigl(\psi_i^{\mathbf{M}}(v_i)\bigr)
\qquad
\text{whenever }\psi_i^{\mathbf{M}}(v_i) \in U_i \cap U_j.
\end{equation}
On points not in $V_i^{\mathrm{reg}}$ or not mapping into $U_i \cap U_j$ we
leave $T_{ij}^{\mathbf{M}}$ undefined; the TSR semantics restrict attention
to in-regime usage as in Definition~\ref{def:tsr}.

\item \textbf{Drift profile.}
We define the drift profile $\varepsilon_{\mathrm{trans}}^{\mathbf{M}}(k)$
for $k$-hop translator paths by composing the one-hop distortions induced by
the chart maps. Informally, for $k=1$ and an in-regime latent entity $u$ with
$z^{\mathbf{M}}(u) \in U_i \cap U_j$ and
$v_i := \chi_i^{\mathbf{M}}(z^{\mathbf{M}}(u))$,
\[
\bigl\|
T_{ij}^{\mathbf{M}}(v_i)
-
\chi_j^{\mathbf{M}}(z^{\mathbf{M}}(u))
\bigr\|
=
0,
\]
and departures from this ideal arise only through observational and encoder
noise (captured in $E_{\mathrm{link}}$) and chart distortion
(Equation~\eqref{eq:chart-distortion}). For general $k$, we define
$\varepsilon_{\mathrm{trans}}^{\mathbf{M}}(k)$ as a non-decreasing function
of $k$ whose leading-order term is a linear accumulation of the one-hop
distortion contributions; Section~6 makes this dependence explicit and shows
that under the bounded-distortion regime of Definition~\ref{def:msr},
$\varepsilon_{\mathrm{trans}}^{\mathbf{M}}(k)$ is first-order in the
distortion profile $\varepsilon_{\mathrm{dist}}^{\mathbf{M}}$.

\end{enumerate}
\end{definition}

By construction, $G_S(\mathbf{M})$ is a TSR of $S$ in the sense of
Definition~\ref{def:tsr}. Intuitively, $G_S$ forgets the manifold and keeps
only the chart-transition structure, turning the Globe back into a
translation graph whose edges represent coordinate changes between overlapping
charts.

\subsection{Constructing \texorpdfstring{$G_S$}{G\_S} on morphisms}
\label{subsec:G-on-morphisms}

We now define $G_S$ on morphisms by discarding the manifold map and retaining
only the feature-space reparameterizations.

\begin{definition}[Functor $G_S$ on morphisms]
\label{def:G-on-morphisms}
Let $H : \mathbf{M} \to \mathbf{M}'$ be a morphism in $\SamGeom^0(S)$ in the
sense of Definition~\ref{def:samgeom-category}, with
\[
H = \bigl(h,\ R\bigr),
\qquad
h : \mathcal{M}^{\mathbf{M}} \to \mathcal{M}^{\mathbf{M}'},
\quad
R = \{R_i : V_i^{\mathbf{M}} \to V_i^{\mathbf{M}'}\}_{i\in\mathcal{I}}.
\]
We define
\[
G_S(H) := R : G_S(\mathbf{M}) \to G_S(\mathbf{M}'),
\]
with tolerances $(\eta_{\mathrm{enc}}, \eta_{\mathrm{trans}})$ to be specified
below. That is, the morphism $G_S(H)$ is simply the family of feature-space
reparameterizations from~$H$, now viewed as a TSR morphism in the sense of
Definition~\ref{def:samtrans-category}.
\end{definition}

The following lemma records that this definition is well-posed.

\begin{lemma}[Well-definedness of $G_S$ on morphisms]
\label{lem:G-morphism}
Let $H = (h,R) : \mathbf{M} \to \mathbf{M}'$ be a morphism in
$\SamGeom^0(S)$ with encoder tolerance $\eta_{\mathrm{enc}}$ and chart
tolerance $\eta_{\mathrm{chart}}$ as in Equations~\eqref{eq:msr-enc-compat}
and~\eqref{eq:msr-chart-compat}. Then
$G_S(H) = R$ is a morphism in $\SamTrans^0(S)$ from
$\mathbf{T}^{\mathbf{M}} = G_S(\mathbf{M})$ to
$\mathbf{T}^{\mathbf{M}'} = G_S(\mathbf{M}')$ with encoder tolerance
$\eta_{\mathrm{enc}}$ and translator tolerance
$\eta_{\mathrm{trans}}$ bounded by a first-order function of
\[
\eta_{\mathrm{chart}},\quad
\varepsilon_{\mathrm{dist}}^{\mathbf{M}}(\cdot),\quad
\varepsilon_{\mathrm{dist}}^{\mathbf{M}'}(\cdot).
\]

\emph{Proof sketch.}
Encoder compatibility for $G_S(H)$ is exactly
Equation~\eqref{eq:msr-enc-compat}, so the encoder tolerance is preserved.

For translators, recall that
$T_{ij}^{\mathbf{M}} = \chi_j^{\mathbf{M}}\circ\psi_i^{\mathbf{M}}$ and
$T_{ij}^{\mathbf{M}'} = \chi_j^{\mathbf{M}'}\circ\psi_i^{\mathbf{M}'}$ on
in-regime domains. Given an in-regime latent entity $u$ and the corresponding
chart coordinates $v_i := \chi_i^{\mathbf{M}}(z^{\mathbf{M}}(u))$ and
$v_i' := \chi_i^{\mathbf{M}'}(z^{\mathbf{M}'}(u))$, the chart compatibility
condition for~$H$ can be written as
\[
\bigl\|R_i(v_i) - R_h^i(v_i)\bigr\|
\le \eta_{\mathrm{chart}},
\qquad
R_h^i := \chi_i^{\mathbf{M}'} \circ h \circ \psi_i^{\mathbf{M}}.
\]
A similar relation holds for $R_j$. The translator compatibility condition we
must check for $G_S(H)$ is
\[
\bigl\|
R_j\bigl(T_{ij}^{\mathbf{M}}(v_i)\bigr)
-
T_{ij}^{\mathbf{M}'}\bigl(R_i(v_i)\bigr)
\bigr\|
\le \eta_{\mathrm{trans}}.
\]
By inserting and subtracting the $h$-induced reparameterizations $R_h^i$ and
$R_h^j$, and then using the distortion bounds
\eqref{eq:chart-distortion} for both $\mathbf{M}$ and $\mathbf{M}'$, one can
show that $\eta_{\mathrm{trans}}$ is controlled by a first-order function of
$\eta_{\mathrm{chart}}$ and the distortion profiles
$\varepsilon_{\mathrm{dist}}^{\mathbf{M}}$, $\varepsilon_{\mathrm{dist}}^{\mathbf{M}'}$.
A detailed proof follows from the distortion bounds and chart-compatibility conditions 
defined above. \qedhere
\end{lemma}

\subsection{Functoriality and gauge stability}
\label{subsec:G-functor}

We now state the functoriality of $G_S$ and its stability with respect to the
error gauges of Definition~\ref{def:error-gauges}.

\begin{proposition}[Functoriality and gauge stability of $G_S$]
\label{prop:G-functor}
The assignment $G_S$ defines a functor
$G_S : \SamGeom^0(S) \to \SamTrans^0(S)$ with the following properties.
\begin{enumerate}[label=(\roman*),leftmargin=*]

\item \textbf{Functoriality.}
For every $\mathbf{M} \in \SamGeom^0(S)$, $G_S(\mathrm{id}_{\mathbf{M}})$ is
the identity morphism on $G_S(\mathbf{M})$, and for any composable morphisms
$H_1 : \mathbf{M}_0 \to \mathbf{M}_1$ and
$H_2 : \mathbf{M}_1 \to \mathbf{M}_2$ in $\SamGeom^0(S)$,
\[
G_S(H_2 \circ H_1) = G_S(H_2) \circ G_S(H_1)
\]
as morphisms in $\SamTrans^0(S)$.

\item \textbf{Gauge stability.}
There exists a first-order slack function
\[
C_G\bigl(
\cdot,\cdot,\cdot
\bigr)
:
\RR_{\ge 0} \times \mathcal{E} \times \mathcal{E} \to \RR_{\ge 0}
\]
(where $\mathcal{E}$ denotes the space of distortion profiles) such that for
any $\mathbf{M},\mathbf{M}' \in \SamGeom^0(S)$,
\begin{equation}
\label{eq:G-gauge-stability}
\Delta_{\SamTrans}\bigl(G_S(\mathbf{M}), G_S(\mathbf{M}')\bigr)
\le
C_G\Bigl(
\Delta_{\SamGeom}(\mathbf{M}, \mathbf{M}'),
\varepsilon_{\mathrm{dist}}^{\mathbf{M}}(\cdot),
\varepsilon_{\mathrm{dist}}^{\mathbf{M}'}(\cdot)
\Bigr),
\end{equation}
with $C_G(0,\varepsilon,\varepsilon)$ first-order in $\varepsilon$.

\end{enumerate}
\end{proposition}

\emph{Proof sketch.}
Functoriality in (i) is immediate from the definition
$G_S(H) = R$: the identity morphism in $\SamGeom^0(S)$ has $h=\mathrm{id}$
and $R_i=\mathrm{id}$ for all $i$, so $G_S(\mathrm{id}_{\mathbf{M}})$ is the
identity morphism on $G_S(\mathbf{M})$. Composition is preserved because
$\bigl(h_2\circ h_1,\{R_i^{(2)}\circ R_i^{(1)}\}\bigr)$ maps under $G_S$ to
$\{R_i^{(2)}\circ R_i^{(1)}\}$, which is exactly the composite of
$G_S(H_1)$ and $G_S(H_2)$ in $\SamTrans^0(S)$.

For (ii), consider two MSRs $\mathbf{M},\mathbf{M}'\in\SamGeom^0(S)$ and the
corresponding TSRs $\mathbf{T}^{\mathbf{M}}$ and
$\mathbf{T}^{\mathbf{M}'}$. The gauge
$\Delta_{\SamTrans}(\mathbf{T}^{\mathbf{M}},\mathbf{T}^{\mathbf{M}'})$ from
Definition~\ref{def:error-gauges} has two terms: the difference between the
drift profiles $\varepsilon_{\mathrm{trans}}^{\mathbf{M}}$ and
$\varepsilon_{\mathrm{trans}}^{\mathbf{M}'}$, and the supremum of translator
disagreements on ideal latent entities. Both are controlled by the chart
structure and the geodesic distances $d_{g^{\mathbf{M}}}$,
$d_{g^{\mathbf{M}'}}$ via Equation~\eqref{eq:chart-distortion}. Using the
definition of $\Delta_{\SamGeom}$, a chaining argument over benign paths,
and the bounded-distortion regime, one can show that the right-hand side of
\eqref{eq:G-gauge-stability} holds with $C_G$ first-order in each of its
arguments. Full details are deferred to Section~6. \qedhere

\subsection{Interpretation and limitations}
\label{subsec:G-discussion}

The functor $G_S$ provides a canonical way to view a well-parameterized MSR
$\mathbf{M}$ as a TSR $\mathbf{T}^{\mathbf{M}}$ whose translators implement
chart transitions. This has two conceptual implications.

First, $G_S$ shows that the translator view and the manifold view are not
competing descriptions but dual coordinatizations of the same semantic
sameness structure $S$. A practitioner who starts from a geometry-first
design (an MSR) can always ``unfold'' it into a translator graph that lies
in $\SamTrans^0(S)$, and error gauges are preserved up to first-order slack
in the distortion profiles.

Second, $G_S$ makes explicit that the stability of the translator view
depends not only on the difference between manifolds
$\Delta_{\SamGeom}(\mathbf{M},\mathbf{M}')$ but also on the intrinsic
distortion of each manifold, as reflected in the dependence of
$C_G$ on $\varepsilon_{\mathrm{dist}}^{\mathbf{M}}$ and
$\varepsilon_{\mathrm{dist}}^{\mathbf{M}'}$ in
\eqref{eq:G-gauge-stability}. This is the geometric analogue of the
translator asymmetry and chart slack terms that appear on the TSR side.

Together with the construction of $F_S$ in
Section~\ref{sec:chartability-F}, the functor $G_S$ completes the bridge
between the translator and manifold design spaces. Section~6 uses these
functors to prove the $\varepsilon$-equivalence theorem
(Theorem~\ref{thm:structural-equivalence}) and, ultimately, the Error Budget
Transfer Theorem, where manifold geometry and error budgets are kept
separate but tightly coupled.
\section{Error-budget transfer and detector-level equivalence}\label{sec:aligned-detectors}
\label{sec:error-transfer}

Sections~\ref{sec:chartability-F} and~\ref{sec:G-construction} constructed
functors
\[
F_S : \SamTrans^0(S) \to \SamGeom^0(S),
\qquad
G_S : \SamGeom^0(S) \to \SamTrans^0(S),
\]
between the well-behaved subcategories of translation-based and manifold-based
realizations of a fixed semantic sameness structure $S$.  The structural
Theorem~\ref{thm:structural-equivalence} showed that these functors form an
$\varepsilon$-equivalence of categories in the sense of
Definition~\ref{def:eps-equivalence}: they preserve realizations up to
first-order slack in the error gauges
$\Delta_{\SamTrans}$ and $\Delta_{\SamGeom}$.

In this section we lift that structural equivalence to the detector level.
We assume that realizations are equipped with sameness detectors in the sense
of Definition~\ref{def:sameness-detector}, and we show that
\emph{error budgets and Davis-style correctness bounds are preserved in
both directions up to first-order slack}.  The key notion is that of
\emph{aligned detectors} on structurally linked realizations.

\subsection{Aligned detectors across realizations}
\label{subsec:aligned-detectors}

Fix a semantic sameness structure $S$ and an in-regime latent subset
$\Omega^{\mathrm{ideal}} \subseteq I$ as in
Definition~\ref{def:error-gauges}.  Let $\mathbf{T}$ be a TSR and
$\mathbf{M}$ an MSR of $S$ that are linked either by chartability
($\mathbf{M} = F_S(\mathbf{T})$) or by geometric realization
($\mathbf{T} = G_S(\mathbf{M})$).  Recall that for each $u \in
\Omega^{\mathrm{ideal}}$ and any modality $i$ with $\pi_i(u)$ defined we
write
\[
v_i^{\mathbf{T}}(u)
  := \varphi_i^{\mathbf{T}}(\pi_i(u)) \in V_i^{\mathbf{T}},
\qquad
z^{\mathbf{M}}(u)
  := \psi_i^{\mathbf{M}}(v_i^{\mathbf{T}}(u)) \in \mathcal{M}^{\mathbf{M}},
\]
which is independent of $i$ on the in-regime set by latent consistency.

\begin{definition}[Aligned detectors]
\label{def:aligned-detectors}
Let $\mathbf{T}$ and $\mathbf{M}$ be as above and let
\[
D^{\mathbf{T}} =
\bigl(
  \mathbf{T},
  \Sigma_{\mathbf{T}}, A_{\mathbf{T}},
  E_{\mathrm{geom}}^{\mathbf{T}},
  E_{\mathrm{link}}^{\mathbf{T}},
  \xi^{\mathbf{T}}, \zeta^{\mathbf{T}},
  \delta_{\mathrm{indep}}^{\mathbf{T}},
  \tau_{\mathrm{vac}}^{\mathbf{T}}
\bigr)
\]
and
\[
D^{\mathbf{M}} =
\bigl(
  \mathbf{M},
  \Sigma_{\mathbf{M}}, A_{\mathbf{M}},
  E_{\mathrm{geom}}^{\mathbf{M}},
  E_{\mathrm{link}}^{\mathbf{M}},
  \xi^{\mathbf{M}}, \zeta^{\mathbf{M}},
  \delta_{\mathrm{indep}}^{\mathbf{M}},
  \tau_{\mathrm{vac}}^{\mathbf{M}}
\bigr)
\]
be sameness detectors on $\mathbf{T}$ and $\mathbf{M}$ in the sense of
Definition~\ref{def:sameness-detector}.  We say that $(D^{\mathbf{T}},
D^{\mathbf{M}})$ are \emph{aligned detectors} for $S$ if:

\begin{enumerate}[label=(\arabic*)]
\item \textbf{Statistic-level alignment.}  There exists
$\eta_{\mathrm{stat}} \ge 0$ such that for all $u \in
\Omega^{\mathrm{ideal}}$ and any modality $i$ with $\pi_i(u)$ defined,
\begin{equation}
\label{eq:stat-alignment}
\Bigl|
  \Sigma_{\mathbf{T}}\bigl(v_i^{\mathbf{T}}(u)\bigr)
  - \Sigma_{\mathbf{M}}\bigl(z^{\mathbf{M}}(u)\bigr)
\Bigr|
\le \eta_{\mathrm{stat}}.
\end{equation}
\item \textbf{Decision-region compatibility.}  The decision rules
$A_{\mathbf{T}}$ and $A_{\mathbf{M}}$ use compatible high-risk thresholds
and abstention regions in the statistic coordinate: there exists
$\eta_{\mathrm{thresh}} \ge 0$ such that the high-risk events
$H^{\mathbf{T}}, H^{\mathbf{M}}$ obey
\[
H^{\mathbf{T}} \triangle H^{\mathbf{M}}
  \subseteq
  \bigl\{ x :
    \bigl|\Sigma_{\mathbf{T}}(x) - \Sigma_{\mathbf{M}}(x)\bigr|
     \le \eta_{\mathrm{thresh}}
  \bigr\}
\]
on non-abstaining, in-regime inputs, where $\triangle$ denotes symmetric
difference.
\item \textbf{Shared calibration and abstention.}  The calibration and
abstention components are either shared or differ only by small slacks:
there exist $\eta_{\mathrm{cal}}, \eta_{\mathrm{abst}} \ge 0$ such that
\[
\bigl|\xi^{\mathbf{T}} - \xi^{\mathbf{M}}\bigr|
\le \eta_{\mathrm{cal}},
\qquad
\bigl|\zeta^{\mathbf{T}} - \zeta^{\mathbf{M}}\bigr|
\le \eta_{\mathrm{abst}}.
\]
\end{enumerate}
We collect the alignment slacks into
\[
\alpha_{\mathrm{align}}
  := \bigl(
      \eta_{\mathrm{stat}},
      \eta_{\mathrm{thresh}},
      \eta_{\mathrm{cal}},
      \eta_{\mathrm{abst}}
    \bigr).
\]
\end{definition}

In practice, aligned detectors arise by design: given a TSR detector
$D^{\mathbf{T}}$ built from features on $\{V_i^{\mathbf{T}}\}$, the
associated manifold detector $D^{\mathbf{M}}$ typically uses the same
parameterization of the statistic on $\mathcal{M}^{\mathbf{M}}$ via
$z^{\mathbf{M}}(u)$ and reuses the same calibration and abstention
policies.  In such cases the alignment slacks are dominated by the
structural chartability parameters and by approximation error in the
statistics.

\subsection{Error-budget equivalence along \texorpdfstring{$F_S$}{F\_S}}
\label{subsec:error-transfer-F}

We first analyze the detector-level effect of the forward functor
$F_S : \SamTrans^0(S) \to \SamGeom^0(S)$.  Recall that for a smoothly
chartable TSR $\mathbf{T}$ the functor constructs a translator metric
quotient and, under the Riemannian realizability assumption, a Davis-style
MSR $\mathbf{M} = F_S(\mathbf{T})$ whose chart structure is compatible
with the translators.  The chartability parameters
\[
\varepsilon_{\mathrm{trans}}^{\mathbf{T}},
\quad
\varepsilon_{\mathrm{dist}}^{\mathbf{M}},
\quad
\delta_{\mathrm{chart}}(\mathbf{T}, \mathbf{M}),
\quad
\delta_{\mathrm{asymm}}(\mathbf{T})
\]
control translator drift, metric distortion, chart overlap slack, and
translator asymmetry, respectively.

We now show that for aligned detectors $(D^{\mathbf{T}}, D^{\mathbf{M}})$
on $(\mathbf{T}, \mathbf{M})$ the error budgets and statistic separation
are equivalent up to first-order slack in these quantities and in the
alignment slacks $\alpha_{\mathrm{align}}$.

\begin{theorem}[Error-budget transfer along $F_S$]
\label{thm:error-transfer-F}
Fix a sameness structure $S$ and a smoothly chartable TSR
$\mathbf{T} \in \SamTrans^0(S)$ with associated MSR
$\mathbf{M} = F_S(\mathbf{T})$.  Let $D^{\mathbf{T}}$ and $D^{\mathbf{M}}$
be aligned detectors in the sense of
Definition~\ref{def:aligned-detectors}.  Write
$E_{\mathrm{geom}}^{\mathbf{R}}$, $E_{\mathrm{link}}^{\mathbf{R}}$,
$\xi^{\mathbf{R}}$, $\zeta^{\mathbf{R}}$, $\delta_{\mathrm{indep}}^{\mathbf{R}}$,
$\Delta_S^{\mathbf{R}}$ for the geometric error, linkage error,
calibration error, abstention failure, independence slack, and statistic
separation of $D^{\mathbf{R}}$ for $\mathbf{R} \in \{\mathbf{T},
\mathbf{M}\}$, with $\Delta_S^{\mathbf{R}}$ defined as in the Davis
framework for the induced statistic.:contentReference[oaicite:0]{index=0}

Then there exist continuous \emph{first-order} slack functions
\[
\Gamma^{F}_{\mathrm{geom}},\;
\Gamma^{F}_{\mathrm{link}},\;
\Gamma^{F}_{\mathrm{sep}},\;
\Gamma^{F}_{\mathrm{det}}
\]
such that:
\begin{align}
\bigl|
  E_{\mathrm{geom}}^{\mathbf{T}} - E_{\mathrm{geom}}^{\mathbf{M}}
\bigr|
&\le
\Gamma^{F}_{\mathrm{geom}}
  \bigl(
    \varepsilon_{\mathrm{trans}}^{\mathbf{T}},
    \varepsilon_{\mathrm{dist}}^{\mathbf{M}},
    \delta_{\mathrm{chart}}(\mathbf{T}, \mathbf{M}),
    \delta_{\mathrm{asymm}}(\mathbf{T}),
    \alpha_{\mathrm{align}}
  \bigr),
\label{eq:F-geom-equivalence}
\\[0.25em]
\bigl|
  E_{\mathrm{link}}^{\mathbf{T}} - E_{\mathrm{link}}^{\mathbf{M}}
\bigr|
&\le
\Gamma^{F}_{\mathrm{link}}
  \bigl(
    \varepsilon_{\mathrm{trans}}^{\mathbf{T}},
    \varepsilon_{\mathrm{dist}}^{\mathbf{M}},
    \delta_{\mathrm{chart}}(\mathbf{T}, \mathbf{M}),
    \delta_{\mathrm{asymm}}(\mathbf{T}),
    \alpha_{\mathrm{align}}
  \bigr),
\label{eq:F-link-equivalence}
\\[0.25em]
\bigl|
  \Delta_S^{\mathbf{T}} - \Delta_S^{\mathbf{M}}
\bigr|
&\le
\Gamma^{F}_{\mathrm{sep}}
  \bigl(
    \varepsilon_{\mathrm{trans}}^{\mathbf{T}},
    \varepsilon_{\mathrm{dist}}^{\mathbf{M}},
    \delta_{\mathrm{chart}}(\mathbf{T}, \mathbf{M}),
    \delta_{\mathrm{asymm}}(\mathbf{T}),
    \alpha_{\mathrm{align}}
  \bigr),
\label{eq:F-sep-equivalence}
\end{align}
and $\Gamma^{F}_{\mathrm{geom}}(0) = \Gamma^{F}_{\mathrm{link}}(0)
= \Gamma^{F}_{\mathrm{sep}}(0) = 0$ with linear (or sublinear) behavior
near~0.

Moreover, let $LB(\cdot)$ denote any Davis-style lower bound on posterior
correctness that is monotone non-decreasing in $\Delta_S$ and monotone
non-increasing in each of $E_{\mathrm{geom}}, E_{\mathrm{link}}, \xi, \zeta$,
and $\delta_{\mathrm{indep}}$; for example, the Cantelli-based bound of
Theorem~2 in the Davis-manifold paper.:contentReference[oaicite:1]{index=1}  Then
\begin{equation}
\label{eq:F-det-equivalence}
\bigl|
  LB(D^{\mathbf{T}}) - LB(D^{\mathbf{M}})
\bigr|
\le
\Gamma^{F}_{\mathrm{det}}
  \bigl(
    \varepsilon_{\mathrm{trans}}^{\mathbf{T}},
    \varepsilon_{\mathrm{dist}}^{\mathbf{M}},
    \delta_{\mathrm{chart}}(\mathbf{T}, \mathbf{M}),
    \delta_{\mathrm{asymm}}(\mathbf{T}),
    \alpha_{\mathrm{align}}
  \bigr),
\end{equation}
with $\Gamma^{F}_{\mathrm{det}}$ first-order in its arguments.
\end{theorem}

\paragraph{Proof sketch.}
The structural $\varepsilon$-equivalence of $F_S$ and the chartability
conditions imply that for in-regime latent entities $u$ the translator
graph and manifold geometry induce path families and distances that agree
up to first-order slack in the parameters
$(\varepsilon_{\mathrm{trans}}^{\mathbf{T}},
\varepsilon_{\mathrm{dist}}^{\mathbf{M}}, \delta_{\mathrm{chart}},
\delta_{\mathrm{asymm}})$.  Aligned statistics then ensure that the
detectors see essentially the same underlying semantic separation along
corresponding paths, and that decision regions differ by at most
$\alpha_{\mathrm{align}}$ in the statistic coordinate.

By definition, $E_{\mathrm{geom}}^{\mathbf{R}}$ and
$E_{\mathrm{link}}^{\mathbf{R}}$ are probabilities of structural and
linkage failures on ideal data.  The chartability bounds give an
inclusion between the corresponding failure events for $\mathbf{T}$ and
$\mathbf{M}$, up to sets of probability controlled by the small
parameters above, yielding the absolute-difference bounds
\eqref{eq:F-geom-equivalence}–\eqref{eq:F-link-equivalence}.  A similar
argument applied to the Davis-style separation parameter $\Delta_S$ gives
\eqref{eq:F-sep-equivalence}.  Finally, monotonicity of $LB(\cdot)$ in
its arguments implies that the induced change in the lower bound is at
most first-order in the changes of
$(E_{\mathrm{geom}}, E_{\mathrm{link}}, \xi, \zeta, \delta_{\mathrm{indep}},
\Delta_S)$, giving \eqref{eq:F-det-equivalence}.  Full details, including
explicit expressions for the slack functions, are deferred to
Section~\ref{sec:structural-equivalence-proofs}.

\subsection{Error-budget equivalence along \texorpdfstring{$G_S$}{G\_S}}
\label{subsec:error-transfer-G}

We now prove the dual result for the geometric-to-translator functor
$G_S : \SamGeom^0(S) \to \SamTrans^0(S)$.  Here the roles of translator
drift and metric distortion are swapped: $\varepsilon_{\mathrm{dist}}$ is
given as part of the MSR, while $\varepsilon_{\mathrm{trans}}$ and
$\delta_{\mathrm{asymm}}$ arise from the induced translators
$T^{\mathbf{M}}_{ij} := \chi^{\mathbf{M}}_j \circ \psi^{\mathbf{M}}_i$.

\begin{theorem}[Error-budget transfer along $G_S$]
\label{thm:error-transfer-G}
Fix a sameness structure $S$ and a geometrically well-parameterized MSR
$\mathbf{M} \in \SamGeom^0(S)$ in the sense of
Definition~\ref{def:geom-well-param}.  Let
$\mathbf{T} := G_S(\mathbf{M})$ be the induced TSR, and let
$D^{\mathbf{M}}$ and $D^{\mathbf{T}}$ be aligned detectors as in
Definition~\ref{def:aligned-detectors}.  Then there exist first-order
slack functions
\[
\Gamma^{G}_{\mathrm{geom}},\;
\Gamma^{G}_{\mathrm{link}},\;
\Gamma^{G}_{\mathrm{sep}},\;
\Gamma^{G}_{\mathrm{det}}
\]
such that:
\begin{align}
\bigl|
  E_{\mathrm{geom}}^{\mathbf{M}} - E_{\mathrm{geom}}^{\mathbf{T}}
\bigr|
&\le
\Gamma^{G}_{\mathrm{geom}}
  \bigl(
    \varepsilon_{\mathrm{dist}}^{\mathbf{M}},
    \varepsilon_{\mathrm{trans}}^{\mathbf{T}},
    \delta_{\mathrm{chart}}(\mathbf{T}, \mathbf{M}),
    \delta_{\mathrm{asymm}}(\mathbf{T}),
    \alpha_{\mathrm{align}}
  \bigr),
\\[0.25em]
\bigl|
  E_{\mathrm{link}}^{\mathbf{M}} - E_{\mathrm{link}}^{\mathbf{T}}
\bigr|
&\le
\Gamma^{G}_{\mathrm{link}}
  \bigl(
    \varepsilon_{\mathrm{dist}}^{\mathbf{M}},
    \varepsilon_{\mathrm{trans}}^{\mathbf{T}},
    \delta_{\mathrm{chart}}(\mathbf{T}, \mathbf{M}),
    \delta_{\mathrm{asymm}}(\mathbf{T}),
    \alpha_{\mathrm{align}}
  \bigr),
\\[0.25em]
\bigl|
  \Delta_S^{\mathbf{M}} - \Delta_S^{\mathbf{T}}
\bigr|
&\le
\Gamma^{G}_{\mathrm{sep}}
  \bigl(
    \varepsilon_{\mathrm{dist}}^{\mathbf{M}},
    \varepsilon_{\mathrm{trans}}^{\mathbf{T}},
    \delta_{\mathrm{chart}}(\mathbf{T}, \mathbf{M}),
    \delta_{\mathrm{asymm}}(\mathbf{T}),
    \alpha_{\mathrm{align}}
  \bigr),
\end{align}
and, for any Davis-style lower bound $LB(\cdot)$ as in
Theorem~\ref{thm:error-transfer-F},
\begin{equation}
\bigl|
  LB(D^{\mathbf{M}}) - LB(D^{\mathbf{T}})
\bigr|
\le
\Gamma^{G}_{\mathrm{det}}
  \bigl(
    \varepsilon_{\mathrm{dist}}^{\mathbf{M}},
    \varepsilon_{\mathrm{trans}}^{\mathbf{T}},
    \delta_{\mathrm{chart}}(\mathbf{T}, \mathbf{M}),
    \delta_{\mathrm{asymm}}(\mathbf{T}),
    \alpha_{\mathrm{align}}
  \bigr),
\end{equation}
with all $\Gamma^{G}$ first-order at the origin.
\end{theorem}

\paragraph{Proof sketch.}
The proof mirrors that of Theorem~\ref{thm:error-transfer-F}.  The
construction of $G_S$ expresses each translator
$T_{ij}^{\mathbf{M}} = \chi^{\mathbf{M}}_j \circ \psi^{\mathbf{M}}_i$ as a
chart-transition map on $\mathcal{M}^{\mathbf{M}}$, so that translator
drift and asymmetry on $\mathbf{T}$ are controlled by the metric
distortion and chart geometry of $\mathbf{M}$.  Geometric well-
parameterization ensures that these quantities are small on the in-regime
subset.  Aligned detectors then again see (up to $\alpha_{\mathrm{align}}$)
the same underlying semantic separation, and the same argument as before
bounds differences in $E_{\mathrm{geom}}$, $E_{\mathrm{link}}$,
$\Delta_S$, and $LB(\cdot)$ by first-order functions of the structural
and alignment parameters.  Details are deferred to
Section~\ref{sec:structural-equivalence-proofs}.

\subsection{Detector-level \texorpdfstring{$\varepsilon$}{ε}-equivalence}
\label{subsec:detector-eps-equiv}

Combining the structural $\varepsilon$-equivalence of
Theorem~\ref{thm:structural-equivalence} with the detector-level transfer
results above yields an $\varepsilon$-equivalence \emph{of detectors} on
$\SamTrans^0(S)$ and $\SamGeom^0(S)$.

\begin{corollary}[Round-trip detector-level equivalence]
\label{cor:detector-equivalence}
Fix $S$ and consider a smoothly chartable TSR
$\mathbf{T} \in \SamTrans^0(S)$ with $\mathbf{M} = F_S(\mathbf{T})$, and a
geometrically well-parameterized MSR
$\mathbf{M}' \in \SamGeom^0(S)$ with $\mathbf{T}' = G_S(\mathbf{M}')$.
Let $D^{\mathbf{T}}$, $D^{\mathbf{M}}$, $D^{\mathbf{M}'}$,
$D^{\mathbf{T}'}$ be aligned detectors on these realizations.

Then there exist first-order slack functions
$C^{F}_{\mathrm{det}}$ and $C^{G}_{\mathrm{det}}$ such that
\begin{align}
\bigl|
  LB(D^{\mathbf{T}}) - LB(D^{F_S(\mathbf{T})})
\bigr|
&\le
C^{F}_{\mathrm{det}}\bigl(\varepsilon_{\mathrm{trans}}^{\mathbf{T}},
\varepsilon_{\mathrm{dist}}^{F_S(\mathbf{T})},
\delta_{\mathrm{chart}}(\mathbf{T}, F_S(\mathbf{T})),
\delta_{\mathrm{asymm}}(\mathbf{T})\bigr),
\\[0.25em]
\bigl|
  LB(D^{\mathbf{M}'}) - LB(D^{G_S(\mathbf{M}')})
\bigr|
&\le
C^{G}_{\mathrm{det}}\bigl(\varepsilon_{\mathrm{dist}}^{\mathbf{M}'},
\varepsilon_{\mathrm{trans}}^{G_S(\mathbf{M}')},
\delta_{\mathrm{chart}}(G_S(\mathbf{M}'), \mathbf{M}'),
\delta_{\mathrm{asymm}}(G_S(\mathbf{M}'))\bigr),
\end{align}
with $C^{F}_{\mathrm{det}}(0) = C^{G}_{\mathrm{det}}(0) = 0$ and first-
order dependence near the origin.

In particular, up to first-order slack in the structural parameters
$(\varepsilon_{\mathrm{trans}}, \varepsilon_{\mathrm{dist}},
\delta_{\mathrm{chart}}, \delta_{\mathrm{asymm}})$, Davis-style correctness
bounds that hold for a manifold-based detector can be transported to a
translator-based detector realizing the same semantic sameness structure,
and conversely.
\end{corollary}

Conceptually, this corollary is the detector-level analogue of the
categorical $\varepsilon$-equivalence in
Theorem~\ref{thm:structural-equivalence}: \emph{within the smoothly
chartable / geometrically well-parameterized regime, working on the
“globe” $\mathbf{M}$ or on a compatible “atlas of flat maps”
$\mathbf{T}$ leads to equivalent guarantees up to first-order slack in
the geometry and chartability parameters}.  This justifies deriving
Davis-style theorems on manifolds and then instantiating them on real
systems implemented as translator graphs, as long as chartability audits
support the small-parameter regime.
\section{Detailed proofs and structural slack bounds}
\label{sec:structural-equivalence-proofs}

This section collects explicit formulas for the slack functions
$\Gamma_{\mathrm{real}}^{F}$, $\Gamma_{\mathrm{real}}^{G}$, $C_F$, and $C_G$,
and proves the structural results announced in Sections~\ref{sec:chartability-F}–\ref{sec:aligned-detectors}.
The central statement is a gauge-stability theorem for the functors
$F_S$ and $G_S$; the error-budget transfer results then appear as corollaries.

Throughout we fix a sameness structure $S$ and work inside the
well–behaved subcategories
$\SamTrans^0(S) \subset \SamTrans(S)$ and $\SamGeom^0(S) \subset \SamGeom(S)$
defined in Sections~\ref{sec:chartability-F}–\ref{sec:G-construction}.
We assume the realizability hypothesis for translator metric quotients
from Section~\ref{sec:chartability-F}.

\subsection{Slack functions and realizability}
\label{subsec:slack-functions}

We first make explicit the ``realizability'' slacks that measure how far a
translator-based or manifold-based realization is from an ideal,
infinitesimally exact geometry, and then define the gauge-stability
slack functions $C_F$ and $C_G$.

\begin{definition}[Realizability slacks]
\label{def:realizability-slacks}
Let $\mathbf{T} \in \SamTrans^0(S)$ be a smoothly chartable TSR with drift
profile $\varepsilon_{\mathrm{trans}}^{\mathbf{T}}(k)$ and asymmetry parameter
$\delta_{\mathrm{asymm}}(\mathbf{T})$ (Definition~\ref{def:translator-asymmetry}).
Let
\[
\tilde{\mathcal{M}}^{\mathbf{T}},\ \tilde d_{\mathbf{T}}
\]
be its translator metric quotient as in
Definition~\ref{def:translator-metric-quotient}, and
let $(\mathcal{M}^{\mathbf{T}}, g^{\mathbf{T}})$ be a Riemannian manifold
and quasi-isometry $q^{\mathbf{T}} : \tilde{\mathcal{M}}^{\mathbf{T}} \to
\mathcal{M}^{\mathbf{T}}$ given by the Riemannian realizability assumption.

We define the \emph{realizability slack} of $\mathbf{T}$ by
\begin{equation}
\label{eq:gamma-real-F}
\Gamma_{\mathrm{real}}^{F}(\mathbf{T})
:= \sup_{[v],[w] \in \tilde{\mathcal{M}}^{\mathbf{T}}}
\left|
\frac{\tilde d_{\mathbf{T}}([v],[w])}{d_{g^{\mathbf{T}}}
\bigl(q^{\mathbf{T}}([v]), q^{\mathbf{T}}([w])\bigr)} - 1
\right|.
\end{equation}
By realizability, $\Gamma_{\mathrm{real}}^{F}(\mathbf{T}) < \infty$ and,
for chartable TSRs, it is controlled by the intrinsic translator
quantities:
\[
\Gamma_{\mathrm{real}}^{F}(\mathbf{T})
\;\lesssim\;
a_1 \sup_{k \ge 1} \varepsilon_{\mathrm{trans}}^{\mathbf{T}}(k)
+ a_2\, \delta_{\mathrm{asymm}}(\mathbf{T}),
\]
for constants $a_1,a_2$ depending only on the ambient architecture and
the chosen path horizon.

Dually, for $\mathbf{M} \in \SamGeom^0(S)$ with chart family
$\{\chi_i^{\mathbf{M}}\}$ and distortion profile
$\varepsilon_{\mathrm{dist}}^{\mathbf{M}}(\cdot)$, we define
\begin{equation}
\label{eq:gamma-real-G}
\Gamma_{\mathrm{real}}^{G}(\mathbf{M})
:= \sup_{z,z' \in \mathcal{M}^{\mathbf{M}}}
\left|
\frac{\delta_i(z,z')}{d_{g^{\mathbf{M}}}(z,z')} - 1
\right|,
\end{equation}
where $\delta_i(z,z') := \|\chi_i^{\mathbf{M}}(z) -
\chi_i^{\mathbf{M}}(z')\|_{V_i}$ on any chart $U_i$ containing $z,z'$.
By Definition~\ref{def:msr-distortion}, this is bounded by the
distortion profile:
\[
\Gamma_{\mathrm{real}}^{G}(\mathbf{M})
\;\lesssim\;
\sup_{\ell > 0} \varepsilon_{\mathrm{dist}}^{\mathbf{M}}(\ell).
\]
Both slacks are \emph{first-order}: whenever the intrinsic translator
errors (for $\mathbf{T}$) or chart distortions (for $\mathbf{M}$) tend
to zero, the corresponding $\Gamma_{\mathrm{real}}^{F}$ or
$\Gamma_{\mathrm{real}}^{G}$ tends to zero.
\end{definition}

We can now define the gauge-level slack functions that appear in the
functorial stability theorem.

\begin{definition}[Gauge-stability slack functions]
\label{def:gauge-slacks}
Let $\Delta_{\SamTrans}$ and $\Delta_{\SamGeom}$ be the error gauges on
$\SamTrans^0(S)$ and $\SamGeom^0(S)$ from
Definition~\ref{def:error-gauges}. We define:
\begin{align}
C_F(\mathbf{T},\mathbf{T}')
&:=
\alpha_F\, \Delta_{\SamTrans}(\mathbf{T},\mathbf{T}')
+ \beta_F
\bigl(
\Gamma_{\mathrm{real}}^{F}(\mathbf{T}) +
\Gamma_{\mathrm{real}}^{F}(\mathbf{T}')
\bigr),
\label{eq:CF-def}
\\[0.3em]
C_G(\mathbf{M},\mathbf{M}')
&:=
\alpha_G\, \Delta_{\SamGeom}(\mathbf{M},\mathbf{M}')
+ \beta_G
\bigl(
\Gamma_{\mathrm{real}}^{G}(\mathbf{M}) +
\Gamma_{\mathrm{real}}^{G}(\mathbf{M}')
\bigr),
\label{eq:CG-def}
\end{align}
for fixed constants $\alpha_F,\beta_F,\alpha_G,\beta_G > 0$ that depend
only on the ambient architectures and on the choice of path horizon and
operational region, not on particular realizations $S$, $\mathbf{T}$, or
$\mathbf{M}$.

By construction these are \emph{first-order} in the gauges and slacks:
if
\[
\Delta_{\SamTrans}(\mathbf{T},\mathbf{T}') \to 0,\quad
\Gamma_{\mathrm{real}}^{F}(\mathbf{T}),\Gamma_{\mathrm{real}}^{F}(\mathbf{T}')
\to 0,
\]
then $C_F(\mathbf{T},\mathbf{T}') \to 0$, and likewise for $C_G$ with
$\Delta_{\SamGeom}$ and $\Gamma_{\mathrm{real}}^{G}$.
\end{definition}

\subsection{Main structural theorem: gauge stability of $F_S$ and $G_S$}
\label{subsec:gauge-stability}

We now prove the main structural result: the functors $F_S$ and $G_S$
are Lipschitz with respect to the error gauges, with slack functions
$C_F$ and $C_G$ from Definition~\ref{def:gauge-slacks}. This theorem is
the ``machine'' that all error-budget transfer statements in
Section~\ref{subsec:aligned-detectors} rest on.

\begin{theorem}[Gauge stability of $F_S$ and $G_S$]
\label{thm:gauge-stability}
Fix a sameness structure $S$ and realizability constants as above.

\begin{enumerate}[]
\item \textbf{Functorial stability of $F_S$.}
For any $\mathbf{T},\mathbf{T}' \in \SamTrans^0(S)$, let
$\mathbf{M} := F_S(\mathbf{T})$ and
$\mathbf{M}' := F_S(\mathbf{T}')$ be the manifold-based realizations
constructed via the translator metric quotient and Riemannian
realizability.
Then
\begin{equation}
\label{eq:gauge-stability-F}
\Delta_{\SamGeom}\bigl(\mathbf{M}, \mathbf{M}'\bigr)
\;\le\;
C_F\bigl(\mathbf{T}, \mathbf{T}'\bigr),
\end{equation}
with $C_F$ as in~\eqref{eq:CF-def}. In particular, if
$\Delta_{\SamTrans}(\mathbf{T},\mathbf{T}') \to 0$ and the realizability
slacks of $\mathbf{T},\mathbf{T}'$ tend to zero, then
$\Delta_{\SamGeom}(\mathbf{M},\mathbf{M}') \to 0$.

\item \textbf{Functorial stability of $G_S$.}
For any $\mathbf{M},\mathbf{M}' \in \SamGeom^0(S)$, let
$\mathbf{T} := G_S(\mathbf{M})$ and
$\mathbf{T}' := G_S(\mathbf{M}')$ be the translator-based realizations
constructed in Section~\ref{sec:G-construction}. Then
\begin{equation}
\label{eq:gauge-stability-G}
\Delta_{\SamTrans}\bigl(\mathbf{T}, \mathbf{T}'\bigr)
\;\le\;
C_G\bigl(\mathbf{M}, \mathbf{M}'\bigr),
\end{equation}
with $C_G$ as in~\eqref{eq:CG-def}. In particular, if
$\Delta_{\SamGeom}(\mathbf{M},\mathbf{M}') \to 0$ and the realizability
slacks of $\mathbf{M},\mathbf{M}'$ tend to zero, then
$\Delta_{\SamTrans}(\mathbf{T},\mathbf{T}') \to 0$.
\end{enumerate}
\end{theorem}

\begin{proof}[Proof sketch]
We sketch the argument for $F_S$; the case of $G_S$ is dual with
translator/metric roles switched.

Let $\mathbf{T},\mathbf{T}' \in \SamTrans^0(S)$, and let
$\mathbf{M}=F_S(\mathbf{T})$,
$\mathbf{M}'=F_S(\mathbf{T}')$ be obtained by:
(i) forming translator metric quotients
$\bigl(\tilde{\mathcal{M}}^{\mathbf{T}}, \tilde d_{\mathbf{T}}\bigr)$ and
$\bigl(\tilde{\mathcal{M}}^{\mathbf{T}'}, \tilde d_{\mathbf{T}'}\bigr)$
as in Definition~\ref{def:translator-metric-quotient};
(ii) invoking realizability to obtain quasi-isometries
$q^{\mathbf{T}} : \tilde{\mathcal{M}}^{\mathbf{T}} \to \mathcal{M}^{\mathbf{T}}$,
$q^{\mathbf{T}'} : \tilde{\mathcal{M}}^{\mathbf{T}'} \to \mathcal{M}^{\mathbf{T}'}$;
and (iii) defining charts $\psi_i^{\mathbf{T}},\psi_i^{\mathbf{T}'}$ by
composing the canonical quotient maps with $q^{\mathbf{T}}, q^{\mathbf{T}'}$.

\emph{Profile part.}
The distortion profiles $\varepsilon_{\mathrm{dist}}^{\mathbf{M}}$ and
$\varepsilon_{\mathrm{dist}}^{\mathbf{M}'}$ compare chart distances
$\delta_i$ to geodesic distances $d_{g^{\mathbf{T}}}$ and
$d_{g^{\mathbf{T}'}}$.
For any pair of latent entities $u,u'$ in the in-regime set, we compare
the corresponding manifold distances via the chain
\[
\left| d_{g^{\mathbf{T}}}(z(u),z(u')) -
       d_{g^{\mathbf{T}'}}(z'(u),z'(u')) \right|
\le A + B + C,
\]
where
\begin{align*}
A &= \left| d_{g^{\mathbf{T}}}(z,z') - \tilde d_{\mathbf{T}}([v],[v']) \right|,
\\
B &= \left| \tilde d_{\mathbf{T}}([v],[v']) - \tilde d_{\mathbf{T}'}([v'],[v'']) \right|,
\\
C &= \left| \tilde d_{\mathbf{T}'}([v'],[v'']) - d_{g^{\mathbf{T}'}}(z',z'') \right|.
\end{align*}
Here $[v],[v'],[v'']$ are the equivalence classes of feature vectors
encoding $u,u'$ in the translator quotients; the exact indexing is not
important.

Terms $A$ and $C$ are controlled by the realizability slacks:
by definition of $\Gamma_{\mathrm{real}}^{F}$,
\[
A \;\lesssim\; \Gamma_{\mathrm{real}}^{F}(\mathbf{T}),\qquad
C \;\lesssim\; \Gamma_{\mathrm{real}}^{F}(\mathbf{T}').
\]
Term $B$ is controlled by the translator gauge
$\Delta_{\SamTrans}(\mathbf{T},\mathbf{T}')$: the quotient metric
$\tilde d_{\mathbf{T}}$ is defined as an infimum over path costs built
from the translators $T_{ij}^{\mathbf{T}}$, and likewise for
$\tilde d_{\mathbf{T}'}$ and $T_{ij}^{\mathbf{T}'}$. By construction of
$\Delta_{\SamTrans}$, the per-hop deviations
$\|T_{ij}^{\mathbf{T}}(v) - T_{ij}^{\mathbf{T}'}(v')\|$ are bounded by
$\Delta_{\SamTrans}(\mathbf{T},\mathbf{T}')$ on in-regime latent points,
and a standard path-comparison argument in metric geometry then yields
\[
B \;\lesssim\; \Delta_{\SamTrans}(\mathbf{T},\mathbf{T}').
\]
Combining the three terms and taking suprema over in-regime latent
pairs shows that the profile component of
$\Delta_{\SamGeom}(\mathbf{M},\mathbf{M}')$ is bounded by a constant
multiple of $C_F(\mathbf{T},\mathbf{T}')$.

\emph{Operational part.}
The operational component of $\Delta_{\SamGeom}$ compares the induced
manifold distances and chart coordinates along benign paths, again over
in-regime latent entities $u$. The same triangle-inequality decomposition
\[
| d_{g^{\mathbf{T}}} - d_{g^{\mathbf{T}'}} |
\;\le\;
| d_{g^{\mathbf{T}}} - \tilde d_{\mathbf{T}} |
+ | \tilde d_{\mathbf{T}} - \tilde d_{\mathbf{T}'} |
+ | \tilde d_{\mathbf{T}'} - d_{g^{\mathbf{T}'}} |
\]
controls these discrepancies uniformly over the operational region, with
the first and last terms again bounded by $\Gamma_{\mathrm{real}}^{F}$
and the middle term by $\Delta_{\SamTrans}$.
This yields the desired bound for the operational piece and thus for
$\Delta_{\SamGeom}(\mathbf{M},\mathbf{M}')$ as a whole.

Collecting constants into $\alpha_F,\beta_F$ produces
\eqref{eq:gauge-stability-F}. The argument for $G_S$ follows the same
template with chart distances and manifold distances swapped and 
$\Gamma_{\mathrm{real}}^{G}$ in place of $\Gamma_{\mathrm{real}}^{F}$,
yielding~\eqref{eq:gauge-stability-G}. \qedhere
\end{proof}

\subsection{Error-budget transfer as corollaries}
\label{subsec:error-budget-corollaries}

We now explain how the main error-budget transfer results in
Section~\ref{sec:aligned-detectors} (Theorems~\ref{cor:EB-F} and
\ref{thm:error-transfer-G}) follow from the gauge-stability theorem
\ref{thm:gauge-stability} together with the Lipschitz dependence of
detector error budgets on the gauges.

Recall that Section~\ref{sec:aligned-detectors} introduces aligned
detectors $D^{\mathbf{R}}$ for $\mathbf{R} \in \SamTrans^0(S)$ or
$\SamGeom^0(S)$, each with an error budget
\[
\bigl(E_{\mathrm{geom}}^{\mathbf{R}},
      E_{\mathrm{link}}^{\mathbf{R}},
      \xi^{\mathbf{R}},
      \zeta^{\mathbf{R}},
      \delta_{\mathrm{indep}}^{\mathbf{R}}\bigr),
\]
and an induced lower bound on correctness for high-risk alerts
$\mathrm{LB}(\mathbf{R})$ of the Davis form.

The key analytic input from Section~\ref{sec:aligned-detectors} is that
these budgets are \emph{Lipschitz in the gauges}: there exist functions
$\Lambda_{\SamTrans}$ and $\Lambda_{\SamGeom}$ such that, for any pair
of realizations $\mathbf{T},\mathbf{T}'$ and $\mathbf{M},\mathbf{M}'$,
\begin{align}
\bigl|
E_{\mathrm{geom}}^{\mathbf{T}} - E_{\mathrm{geom}}^{\mathbf{T}'}
\bigr|
&\le \Lambda_{\SamTrans}\bigl(
\Delta_{\SamTrans}(\mathbf{T},\mathbf{T}')
\bigr),
\quad\text{and likewise for } E_{\mathrm{link}},\xi,\zeta,\delta_{\mathrm{indep}},
\label{eq:EB-Lip-trans}
\\
\bigl|
E_{\mathrm{geom}}^{\mathbf{M}} - E_{\mathrm{geom}}^{\mathbf{M}'}
\bigr|
&\le \Lambda_{\SamGeom}\bigl(
\Delta_{\SamGeom}(\mathbf{M},\mathbf{M}')
\bigr),
\label{eq:EB-Lip-geom}
\end{align}
with $\Lambda_{\SamTrans}$, $\Lambda_{\SamGeom}$ first-order at the
origin.

We now combine these Lipschitz properties with
Theorem~\ref{thm:gauge-stability}.

\begin{corollary}[Error-budget equivalence along $F_S$]
\label{cor:EB-F}
Let $\mathbf{T},\mathbf{T}' \in \SamTrans^0(S)$ and
$\mathbf{M} = F_S(\mathbf{T})$,
$\mathbf{M}' = F_S(\mathbf{T}')$. Let $D^{\mathbf{T}}$ and
$D^{\mathbf{T}'}$ be aligned detectors on $\mathbf{T}$ and $\mathbf{T}'$
with Davis-style error budgets and lower bounds $\mathrm{LB}(\mathbf{T})$,
$\mathrm{LB}(\mathbf{T}')$, and let $D^{\mathbf{M}}$,
$D^{\mathbf{M}'}$ be their aligned MSR detectors with bounds
$\mathrm{LB}(\mathbf{M})$, $\mathrm{LB}(\mathbf{M}')$.

Then there exists a first-order function
$\Gamma_F^{\mathrm{det}}$ (depending only on the Lipschitz moduli in
\eqref{eq:EB-Lip-trans}–\eqref{eq:EB-Lip-geom} and on the constants in
$C_F$) such that
\[
\bigl|
\mathrm{LB}(\mathbf{M}) - \mathrm{LB}(\mathbf{M}')
\bigr|
\;\le\;
\Gamma_F^{\mathrm{det}}\bigl(
\Delta_{\SamTrans}(\mathbf{T},\mathbf{T}'),
\Gamma_{\mathrm{real}}^{F}(\mathbf{T}),
\Gamma_{\mathrm{real}}^{F}(\mathbf{T}')
\bigr),
\]
and similarly with $(\mathbf{M},\mathbf{M}')$ and
$(\mathbf{T},\mathbf{T}')$ reversed:
\[
\bigl|
\mathrm{LB}(\mathbf{T}) - \mathrm{LB}(\mathbf{T}')
\bigr|
\;\le\;
\Gamma_F^{\mathrm{det}}\bigl(
\Delta_{\SamTrans}(\mathbf{T},\mathbf{T}'),
\Gamma_{\mathrm{real}}^{F}(\mathbf{T}),
\Gamma_{\mathrm{real}}^{F}(\mathbf{T}')
\bigr).
\]
In particular, if $\Delta_{\SamTrans}(\mathbf{T},\mathbf{T}')$ and the
realizability slacks of $\mathbf{T},\mathbf{T}'$ tend to zero, then all
four lower bounds
$\mathrm{LB}(\mathbf{T}),\mathrm{LB}(\mathbf{T}'),
 \mathrm{LB}(\mathbf{M}),\mathrm{LB}(\mathbf{M}')$
converge together.
\end{corollary}

\begin{proof}
Apply the Lipschitz property~\eqref{eq:EB-Lip-geom} to the pair
$(\mathbf{M},\mathbf{M}')$ and then use the gauge-stability bound
\eqref{eq:gauge-stability-F}:
\[
\bigl|
\mathrm{LB}(\mathbf{M}) - \mathrm{LB}(\mathbf{M}')
\bigr|
\;\lesssim\;
\Lambda_{\SamGeom}\bigl(
\Delta_{\SamGeom}(\mathbf{M},\mathbf{M}')
\bigr)
\;\le\;
\Lambda_{\SamGeom}\bigl(
C_F(\mathbf{T},\mathbf{T}')
\bigr).
\]
Since $C_F$ is first-order in
$\Delta_{\SamTrans}$ and $\Gamma_{\mathrm{real}}^{F}$, and
$\Lambda_{\SamGeom}$ is first-order at the origin, the composite
$\Gamma_F^{\mathrm{det}} := \Lambda_{\SamGeom} \circ C_F$ has the
claimed first-order dependence.

The same reasoning applied directly to
$(\mathbf{T},\mathbf{T}')$ via~\eqref{eq:EB-Lip-trans} bounds
$|\mathrm{LB}(\mathbf{T}) - \mathrm{LB}(\mathbf{T}')|$ by a (possibly
different but comparable) first-order function of
$\Delta_{\SamTrans}$ and $\Gamma_{\mathrm{real}}^{F}$; we absorb both
into the single $\Gamma_F^{\mathrm{det}}$.
\end{proof}

\begin{corollary}[Error-budget equivalence along $G_S$]
\label{cor:EB-G}
Let $\mathbf{M},\mathbf{M}' \in \SamGeom^0(S)$ and
$\mathbf{T} = G_S(\mathbf{M})$,
$\mathbf{T}' = G_S(\mathbf{M}')$. With aligned detectors and Davis-style
error budgets as in Corollary~\ref{cor:EB-F}, there exists a first-order
function $\Gamma_G^{\mathrm{det}}$ such that
\[
\bigl|
\mathrm{LB}(\mathbf{M}) - \mathrm{LB}(\mathbf{M}')
\bigr|,
\quad
\bigl|
\mathrm{LB}(\mathbf{T}) - \mathrm{LB}(\mathbf{T}')
\bigr|
\;\le\;
\Gamma_G^{\mathrm{det}}\bigl(
\Delta_{\SamGeom}(\mathbf{M},\mathbf{M}'),
\Gamma_{\mathrm{real}}^{G}(\mathbf{M}),
\Gamma_{\mathrm{real}}^{G}(\mathbf{M}')
\bigr).
\]
In particular, if $\Delta_{\SamGeom}(\mathbf{M},\mathbf{M}')$ and the
realizability slacks of $\mathbf{M},\mathbf{M}'$ tend to zero, all four
lower bounds converge together.
\end{corollary}

\begin{proof}
Identical in structure to Corollary~\ref{cor:EB-F}, using
\eqref{eq:gauge-stability-G} and the Lipschitz property
\eqref{eq:EB-Lip-trans} for TSR error budgets.
\end{proof}

Together, Theorem~\ref{thm:gauge-stability} and
Corollaries~\ref{cor:EB-F}–\ref{cor:EB-G} formalize the idea that the
TSR and MSR descriptions of a given sameness structure $S$ are
\emph{equivalent as detection substrates}: small gauge distance between
two TSRs implies small gauge distance between their manifold
realizations, and vice versa, and aligned detectors on either side
inherit this equivalence at the level of error budgets and lower
correctness bounds.
\section{Operational audits and evaluation protocols}
\label{sec:audits}

This section specifies empirical audits and evaluation protocols for Davis systems.
It is deliberately \emph{procedural} and does not report numerical results, case
studies, or deployment performance. The goal is to make the theoretical assumptions
of Sections~3–5 falsifiable in practice: each audit targets a specific part of the
geometry, separation, or error budget, and is designed to be run on held-out data
or prospective streams.

Throughout, we fix a Davis system
\[
\mathcal{D}
= \bigl( (M,g,\delta), P(L), c, h, S, A, \text{error budget}\bigr)
\]
in some domain, with distortion profile $\varepsilon(L)$, configuration margins
$(\kappa_{\mathrm{hard}},\kappa_{\mathrm{soft}})$, and error-budget components
$(E_{\mathrm{geom}}, E_{\mathrm{link}}, \xi, \zeta, \delta_{\mathrm{indep}})$
as defined earlier. All protocols assume strict train/validation splits and
time-asymmetric estimation where applicable (no look-ahead).

\subsection{Evaluation principles and pre-registration}

Before running any audit or reporting any number, we recommend pre-registering:

\begin{enumerate}[label=(\roman*)]
\item \textbf{Primary endpoints.} Which quantities will be assessed? Examples:
  (a) distortion tails $\varepsilon(L_\star)$ at the proposed operational horizon,
  (b) effective separation gap $\kappa_{\mathrm{soft}} - 2R\varepsilon(L_\star)$,
  (c) empirical Cantelli bound tightness, (d) calibration error $\xi$ in a high-risk
  band, (e) abstention coverage and residual failure $\zeta$.
\item \textbf{Path and slice selection.} How will benign paths in $P(L)$ and
  evaluation slices be sampled (e.g., by time, region, demographic group,
  capture condition)? Which slices count as in-distribution $R_{\mathrm{valid}}$?
\item \textbf{Thresholds and vacuity criteria.} Pre-specify acceptable ranges
  for distortion (e.g., $\varepsilon(L_\star) \le \varepsilon^\star$), error-budget
  sums $E_{\mathrm{geom}} + E_{\mathrm{link}} + \xi + \zeta \le \tau_{\mathrm{vac}}$,
  and independence slack $\delta_{\mathrm{indep}}$.
\item \textbf{Baselines and ablations.} Define baselines (e.g., flat Euclidean
  models, models without smoothness regularization, non-geometric detectors),
  and which components will be ablated (e.g., remove path features,
  remove auxiliary detectors, remove abstention).
\item \textbf{Data governance and leakage.} Fix time windows for training vs\
  evaluation, forbid re-using evaluation data for tuning, and record configuration
  hashes for all runs to support auditability.
\end{enumerate}

All audits below should be run on held-out validation data (retrospective or
prospective) that is not used to train $f_\theta$, fit $S$, or tune the abstention
policy $A$.

\subsection{Distortion audits: verifying the bounded-distortion regime}

These audits target the geometric assumptions behind the distortion profile
$\varepsilon(L)$ and the path-horizon choice $L_\star$ (Theorem~1 and the
non-vacuity condition $\kappa_{\mathrm{soft}} - 2R\varepsilon(L_\star) > 0$).

\paragraph{Protocol 1: geodesic–Euclidean distortion vs.~path length.}

\begin{enumerate}[label=(\arabic*)]
\item \textbf{Sample benign paths.} On validation data, construct a set of
  identity-preserving paths $\{\gamma_m\}_{m=1}^M \subset P(L_{\max})$ via
  natural temporal sequences (e.g., time windows, short trajectories) or
  controlled perturbations. Record their lengths $L(\gamma_m)$.
\item \textbf{Subdivide into segments.} For each $\gamma_m$, sample pairs
  of times $0 \le s < t \le 1$ with arc-length distance
  $d_g\bigl(z(\gamma_m(s)), z(\gamma_m(t))\bigr)$ spanning a grid of target
  lengths $\ell \in \{\ell_1,\dots,\ell_K\}$ up to $L_{\max}$.
\item \textbf{Approximate geodesic distance.} For each pair $(s,t)$, approximate
  $d_g$ using either:
  \begin{itemize}
  \item the path length of the interpolated trajectory in the embedding
        (if $g$ is induced by a pullback metric), or
  \item a discrete geodesic approximation (e.g., graph-based shortest paths)
        constrained to lie within a local neighborhood.
  \end{itemize}
\item \textbf{Compute distortion ratios.} Let $\delta$ be the ambient
  chord metric (Euclidean distance in the embedding). For each pair, compute
  the distortion ratio
  \[
  \rho(s,t)
  = \frac{\delta\bigl(z(\gamma_m(s)), z(\gamma_m(t))\bigr)}
         {d_g\bigl(z(\gamma_m(s)), z(\gamma_m(t))\bigr)}.
  \]
  Aggregate $\rho$ into bins by geodesic length $\ell$.
\item \textbf{Summarize distortion.} For each length bin $\ell$, estimate
  $\mathbb{E}[|\rho-1|\mid d_g \approx \ell]$ and upper quantiles
  (e.g., 95th/99th percentile). Plot or tabulate these as a ``distortion curve''
  $\widehat{\varepsilon}_{\mathrm{geom}}(\ell)$.
\item \textbf{Choose $L_\star$ and validate.} Select an operational horizon
  $L_\star$ such that $\widehat{\varepsilon}_{\mathrm{geom}}(\ell) \le \varepsilon^\star$
  for all $\ell \le L_\star$, with $\varepsilon^\star$ chosen so that the effective
  soft margin $\kappa_{\mathrm{soft}} - 2R\varepsilon^\star$ remains positive.
  Document any slices (e.g., specific conditions or subpopulations) where
  distortion tails are heavy; these may be excluded from $R_{\mathrm{valid}}$
  or treated as out-of-distribution.
\end{enumerate}

\paragraph{Protocol 2: curvature and smoothness along paths.}

To empirically test the smoothness assumptions used in Section~3:

\begin{enumerate}[label=(\arabic*)]
\item Approximate discrete curvature along each path by the angles between
  successive segments in the embedding, or by second differences of $z(t)$
  in arc-length parameterization.
\item Confirm that curvature statistics (mean, 95th percentile) remain below
  a pre-registered threshold for paths with length $\le L_\star$.
\item Flag regimes where curvature spikes (e.g., hard scene cuts, recombination
  events, abrupt control changes) as candidates for abstention or separate
  modeling.
\end{enumerate}

These two audits jointly support (T1) in Section~6.2: they test whether the
bounded-distortion regime is empirically valid on the intended deployment
distribution.

\subsection{Configuration and margin audits}

These audits target the configuration map $c$, coarse labels $h$, and the
separation margins $(\kappa_{\mathrm{hard}},\kappa_{\mathrm{soft}})$ used in
the Cantelli bound and non-vacuity conditions.

\paragraph{Protocol 3: separation in geodesic and ambient distance.}

\begin{enumerate}[label=(\arabic*)]
\item \textbf{Collect labeled pairs.} On validation data, assemble a set of
  endpoint pairs $(x_a,x_b)$ with coarse labels
  $h\bigl(c(z(x_a))\bigr)$ and $h\bigl(c(z(x_b))\bigr)$, distinguishing
  ``similar'' vs.\ ``changed'' vs.\ ``ambiguous'' configurations.
\item \textbf{Measure distances.} For each pair, compute both the approximate
  geodesic distance $d_g\bigl(z(x_a),z(x_b)\bigr)$ and ambient distance
  $\delta\bigl(z(x_a),z(x_b)\bigr)$.
\item \textbf{Estimate margins.} For each label pair (e.g., similar–similar,
  similar–changed, changed–changed), estimate the empirical distribution of
  distances. From these, estimate:
  \[
  \widehat{\kappa}_{\mathrm{hard}}
  \approx \inf_{(a,b)\ \text{cross-config}} d_g(z_a,z_b),\qquad
  \widehat{\kappa}_{\mathrm{soft}}
  \approx \mathrm{quantile}_{q}(d_g(z_a,z_b))
  \]
  for a pre-registered quantile $q$ (e.g., 5th percentile of cross-config pairs).
\item \textbf{Check non-vacuity under distortion.} Using the distortion
  estimates from the previous subsection, verify that
  $\widehat{\kappa}_{\mathrm{soft}} - 2R\widehat{\varepsilon}_{\mathrm{geom}}(L_\star) > 0$
  on the slices of interest. If this inequality fails or is marginal,
  treat the corresponding slice as ambiguous and rely on abstention
  or auxiliary detectors.
\end{enumerate}

\paragraph{Protocol 4: stability across slices.}

Repeat Protocol~3 across pre-registered slices (e.g., demographics, capture
conditions, time periods) and compare the estimated margins; large variation
may indicate bias in $f_\theta$ or in the configuration definition. Where
gaps appear, consider slice-specific abstention policies or separate retraining.

\subsection{Feature--geometry linkage audits}

These audits target the monotone link $\Delta S \ge g(\Delta g)$ between
configuration movement in geometry and changes in the scalar statistic $S$.

\paragraph{Protocol 5: monotone linkage in bins of geodesic displacement.}

\begin{enumerate}[label=(\arabic*)]
\item \textbf{Define path segments and labels.} For each benign path
  $\gamma \in P(L_\star)$ and time pair $(s,t)$, compute the geodesic
  displacement $\Delta g = d_g\bigl(z(\gamma(s)),z(\gamma(t))\bigr)$ and
  the corresponding feature-based change in the statistic, e.g.,
  $\Delta S = S_t - S_s$ or $S_t$ alone if $S$ is cumulative.
\item \textbf{Bin by geodesic distance.} Partition the range of $\Delta g$
  into bins $B_1,\dots,B_K$ on an operational interval
  $[0,g_{\max}]$. For each bin, compute:
  \[
  \overline{\Delta S}_k
  = \mathbb{E}[\Delta S \mid \Delta g \in B_k],\qquad
  \widehat{\mathrm{Var}}(\Delta S \mid \Delta g \in B_k).
  \]
\item \textbf{Test monotonicity.} Fit a monotone regression
  $\widehat{g}$ of $\overline{\Delta S}_k$ on the midpoints of $B_k$,
  and measure deviations from monotonicity (e.g., number and magnitude
  of violations). Pre-register an acceptable tolerance for such violations.
\item \textbf{Report linkage slack.} Summarize a linkage slack term
  $\widehat{\eta}_{\mathrm{link}}$ capturing how far the empirical relation
  deviates from an ideal monotone link; this feeds into the linkage
  component of the error budget $E_{\mathrm{link}}$.
\end{enumerate}

\paragraph{Protocol 6: variance and finite-moment checks.}

To justify finite-variance Cantelli bounds, estimate second moments of
$S$ (or $\Delta S$) conditional on configuration changes. If heavy tails
are observed, tighten the operational range or incorporate robust statistics
(e.g., winsorized $S$) and re-audit.

\subsection{Calibration, abstention, and error-budget estimation}

These audits estimate the non-geometric terms $(\xi,\zeta)$ and link
theoretical error budgets to empirical behavior.

\paragraph{Protocol 7: calibration in the high-risk region.}

\begin{enumerate}[label=(\arabic*)]
\item \textbf{Fit and freeze calibration.} Using historical training data,
  fit a monotone calibration map $\widehat{\pi} = g_{\mathrm{cal}}(S)$ from
  the statistic to an interpretable risk quantity (e.g., probability of
  class change within horizon $T$). Freeze $g_{\mathrm{cal}}$ before
  calibration auditing.
\item \textbf{Evaluate on validation.} On held-out or prospective data,
  compute calibration metrics (Brier score, expected calibration error,
  reliability curves) on:
  \begin{itemize}
  \item the full support of $S$, and
  \item a high-risk band (e.g., $S$ above a pre-registered threshold where
        actions are taken).
  \end{itemize}
\item \textbf{Estimate $\xi$.} Define $\widehat{\xi}$ as the maximum
  deviation between empirical and nominal probabilities in the high-risk
  band, or as a pre-registered functional of the calibration curves.
  This quantity enters the error budget as the calibration term.
\end{enumerate}

\paragraph{Protocol 8: abstention behavior and $\zeta$.}

\begin{enumerate}[label=(\arabic*)]
\item \textbf{Record abstentions.} Under a candidate abstention policy $A$,
  log for each evaluation sample whether the system abstained and the
  rationale (e.g., large distortion, high OOD score, large predictive
  interval).
\item \textbf{Estimate coverage.} For each slice, estimate coverage
  $1 - A$ (fraction of non-abstained cases) and the conditional error
  rates among non-abstained decisions (e.g., misclassification rate or
  violation of a domain-specific safety criterion).
\item \textbf{Estimate $\zeta$.} Define $\widehat{\zeta}$ as the frequency
  with which the system \emph{fails} to abstain in regimes where distortion
  or OOD scores exceed pre-registered safety thresholds. This captures
  abstention failures in the error budget.
\end{enumerate}

\paragraph{Protocol 9: empirical error-budget decomposition.}

On validation data where ground-truth outcomes are available, label each
failure with its dominant source (geometry drift, linkage failure, calibration
error, abstention failure, or uncontrollable noise). Estimate empirical
frequencies of each error type and compare to the theoretical components
$(E_{\mathrm{geom}},E_{\mathrm{link}},\xi,\zeta)$; large discrepancies signal
either model misspecification or unmodeled dependencies.

\subsection{Independence slack and compositional bounds}

The compositional dominance bound in Theorem~2 assumes approximate
conditional independence of the error sources; when this fails, the theory
recommends a conservative union bound. This subsection sketches how to
empirically assess the independence slack $\delta_{\mathrm{indep}}$.

\paragraph{Protocol 10: independence diagnostics.}

\begin{enumerate}[label=(\arabic*)]
\item \textbf{Error indicators.} For each sample, define Bernoulli indicators
  $E_{\mathrm{geom}}, E_{\mathrm{link}}, E_{\mathrm{cal}}, E_{\mathrm{abst}}$
  for the occurrence of each error type, based on the audits above and
  domain-specific thresholds.
\item \textbf{Estimate joint vs.\ product.} For selected pairs or triples
  of error types (e.g., $(E_{\mathrm{geom}},E_{\mathrm{link}})$,
  $(E_{\mathrm{link}},E_{\mathrm{cal}})$), estimate:
  \[
  p_{12} = \mathbb{P}(E_1 \land E_2),\qquad
  p_1 = \mathbb{P}(E_1),\quad p_2 = \mathbb{P}(E_2),
  \]
  and the deviation from independence
  $\Delta_{12} = p_{12} - p_1 p_2$.
\item \textbf{Define $\delta_{\mathrm{indep}}$.} Aggregate the magnitudes
  $|\Delta_{12}|$ across pairs into an independence slack
  $\widehat{\delta}_{\mathrm{indep}}$ (e.g., maximum or suitable norm),
  and compare the compositional bound using products to the conservative
  union bound using sums. If $\widehat{\delta}_{\mathrm{indep}}$ is large,
  prefer the union bound in practice.
\end{enumerate}

\subsection{Slicing, OOD monitoring, and lifecycle audits}

Finally, we outline cross-cutting audits that track geometry and performance
over time and across slices.

\paragraph{Protocol 11: slice-wise Davis audits.}

For each pre-registered slice (e.g., region, demographic group, capture
condition, hardware profile):

\begin{enumerate}[label=(\arabic*)]
\item Re-run Protocols~1–9 restricted to that slice.
\item Compare distortion curves, margins, linkage slack, calibration error,
  and abstention behavior across slices.
\item Flag slices where any component of the error budget becomes vacuous
  (e.g., $E_{\mathrm{geom}} + E_{\mathrm{link}} + \xi + \zeta > \tau_{\mathrm{vac}}$),
  and treat them as out-of-regime for the current system.
\end{enumerate}

\paragraph{Protocol 12: OOD detection and retraining triggers.}

Monitor distributions of:
\begin{itemize}
\item distortion ratios $\rho(s,t)$,
\item path curvature statistics,
\item feature distributions used in $S$,
\item OOD or density scores from auxiliary detectors.
\end{itemize}
If these drift beyond pre-registered control limits, increase abstention,
freeze thresholds, and consider retraining $f_\theta$ or redefining
$P(L)$ and $L_\star$ before resuming full operation.

\paragraph{Section scope note.}

This section specifies \emph{audit protocols and reporting templates} for
future empirical work with Davis systems. By design, it contains no numerical
results, plots, or case studies. In practical deployments, we recommend that
these protocols be instantiated as part of a comprehensive system dossier,
with results refreshed on a fixed cadence and after major system changes.
\section{Case study: a geometry-first multi-modal threat detector}
\label{sec:kraken-case-study}

The abstract framework of Sections~\ref{sec:sameness}--\ref{sec:structural-equivalence-proofs}
was developed independently of any particular domain. In parallel, the author built
a deployed system---codenamed \emph{KRAKEN} in a separate, application-focused
manuscript---for real-time, multi-modal maritime threat detection.%
\footnote{The present section intentionally omits operational details (sensor
layouts, campaign locations, adversary tactics) and focuses only on the structural
aspects relevant to the TSR--MSR theory: semantic sameness, translator graphs,
geometric realizations, and error budgets.}
In hindsight, KRAKEN can be understood as a concrete instance of the TSR--MSR
equivalence and error-budget decomposition developed in this paper.

This section sketches that correspondence. It is not required for the theoretical
results, but serves as a sanity check: a complex, safety-critical system built
before the formalism was articulated nonetheless fits naturally into the
semantic-sameness, realization, and functorial picture.

\subsection{Operational setting (informal)}

At a high level, the system fuses heterogeneous sensor streams to detect and track
subsurface threats and undersea infrastructure hazards in coastal waters.
Representative modalities include:
\begin{itemize}
  \item \textbf{Kinetic} channels: distributed acoustic sensing (fiber strain),
        passive SONAR spectrograms, and related vibration measurements;
  \item \textbf{Electromagnetic} channels: synthetic aperture radar and optical/IR
        imagery of the sea surface, plus RF emissions from platforms and shore;
  \item \textbf{Contextual} channels: shipping lanes, bathymetry, environmental
        metadata, and operator annotations.
\end{itemize}
Each modality has its own units, SNR regime, failure modes, and coverage gaps.
Operators care about \emph{entities}: specific submarines, vessels, and patterns
of behavior. The core question is semantic: when is a faint acoustic trace,
a marginal image artifact, and a weak RF transient evidence of the \emph{same}
latent threat, and when are they unrelated clutter?

\subsection{Semantic sameness structure for maritime threats}

We briefly sketch a semantic sameness structure
$S_{\mathrm{mar}} = (I, \{X_i\}_{i\in\mathcal{I}}, \{\pi_i\}_{i\in\mathcal{I}},
\approx, \mathcal{P}_S(L))$ in the sense of Section~\ref{sec:sameness} that
captures this setting.

\begin{itemize}
  \item \textbf{Latent space $I$.} Elements $u \in I$ represent latent maritime
  \emph{threat states} at an instant: a tuple $(C,E,Q)$ of threat class
  (platform or behavior type), environment (sea state, sound-speed profile,
  background traffic), and configuration (speed, depth, machinery state, loadout).

  \item \textbf{Modalities $\{X_i\}$.} Each $i\in\mathcal{I}$ indexes a sensor
  family (e.g., DAS fiber segments, SONAR beams, satellite tiles, RF channels).
  $X_i$ is the space of idealized observations from that family over a short
  time window.

  \item \textbf{Rendering maps $\pi_i : I \to X_i$.} Given a latent state
  $u = (C,E,Q)$, the map $\pi_i(u)$ encodes what an ideal, noise-free sensor
  of modality $i$ \emph{would} observe: an acoustic spectrum induced by
  propeller cavitation, a surface wake pattern in SAR, a strain pattern along
  a buried fiber, an RF emission mask, and so on. In practice, real observations
  include noise, occlusion, and quantization; as in Section~\ref{sec:sameness},
  we blur notation and treat $x_i\in X_i$ as an observation ``of'' $u$ when it
  lies within an acceptable noise ball around $\pi_i(u)$.

  \item \textbf{Sameness relation $\approx$.} Two observations $x\in X_i$,
  $x'\in X_j$ are semantically equivalent, $x\approx x'$, if they are rendered
  from the \emph{same} latent entity within a short temporal tolerance: there
  exists a path $\gamma\in\mathcal{P}_S(L)$ and time $t$ such that
  $x\approx \pi_i(\gamma(t))$, $x'\approx \pi_j(\gamma(t))$.
  Intuitively, they are evidence about the same contact rather than coincident
  clutter.

  \item \textbf{Benign path family $\mathcal{P}_S(L)$.} Paths
  $\gamma : [0,1]\to I$ describe physically plausible threat trajectories over
  an operational horizon $L$: slow drifts in position and depth, machinery
  state changes, and routine environmental variation. This family encodes both
  how fast threats can move in latent space and which directions are benign.
\end{itemize}

KRAKEN was designed without this formal vocabulary, but its informal threat
models and scenario design align closely with $S_{\mathrm{mar}}$.

\subsection{TSR instantiation: translator graph over sensor embeddings}

On the translator side, KRAKEN provides a TSR
$\mathbf{T}_{\mathrm{mar}} \in \SamTrans^0(S_{\mathrm{mar}})$ in the sense of
Definition~\ref{def:tsr}.

\begin{itemize}
  \item \textbf{Feature spaces $V_i$ and encoders $\varphi_i$.} Each modality
  has an encoder $\varphi_i : X_i \to V_i$ (CNNs for spectrograms, transformers
  for sequences, etc.) producing modality-specific embeddings. These are the
  $V_i$ in $\mathbf{T}_{\mathrm{mar}}$.

  \item \textbf{Translator graph.} A sparse, directed graph $G_T$ connects
  modalities via learned translators $T_{ij} : V_i \to V_j$.
  Edges exist where physics and coverage justify translation (e.g., DAS$\to$SONAR,
  SONAR$\to$DAS, DAS$\to$SAT, RF$\to$SAT). Translators are trained and calibrated
  to satisfy per-edge error bounds
  \[
    \bigl\| T_{ij}(v_i) - v_j \bigr\| \le \varepsilon_{ij}
  \]
  on in-regime data, yielding a drift profile $\varepsilon_{\mathrm{trans}}(k)$
  for $k$-hop paths as in Definition~\ref{def:tsr}. Spectral routing on $G_T$
  avoids paths whose accumulated $\varepsilon_{\mathrm{trans}}$ exceeds a
  configured budget.

  \item \textbf{Chartability side conditions.} Empirically, the learned
  translators on the operational regime satisfy:
  \begin{itemize}
    \item approximate local bi-Lipschitz behavior (bounded Jacobians);
    \item small cocycle residual
      $\|T_{jk}(T_{ij}(v)) - T_{ik}(v)\|$ on triple overlaps;
    \item bounded asymmetry $\|T_{ji}(T_{ij}(v)) - v\|$ on round-trips.
  \end{itemize}
  These are precisely the smooth-chartability conditions from
  Section~\ref{sec:chartability-F}, verified by held-out audits rather than
  assumed.
\end{itemize}

From the standpoint of this paper, KRAKEN’s representation layer is a richly
engineered TSR with a calibrated error budget and empirical evidence of smooth
chartability.

\subsection{MSR instantiation: a Davis-style threat manifold}

On the geometric side, KRAKEN instantiates an MSR
$\mathbf{M}_{\mathrm{mar}} \in \SamGeom^0(S_{\mathrm{mar}})$ in the spirit of
Definition~\ref{def:msr} and the Davis-manifold construction.

\begin{itemize}
  \item \textbf{Threat manifold $(\mathcal{M}, g)$.} Multi-modal evidence fused
  through the translator graph is encoded into a low-dimensional Riemannian
  manifold $(\mathcal{M}, g)$ using a HERALD-style encoder trained with a
  contrastive objective and a smoothness regularizer. Geodesic distance $d_g$
  aligns with threat semantics: quiet and loud variants of the same platform
  lie close; distinct classes separate with a soft margin.

  \item \textbf{Charts and chart distances.} For each modality (or fused
  representation) there is a chart map
  $\chi_i : U_i \subset \mathcal{M} \to V_i$ with bounded distortion in the
  sense of Definition~\ref{def:msr}. Distortion audits along benign paths
  in $\mathcal{P}_S(L)$ empirically validate a small distortion profile
  $\varepsilon_{\mathrm{dist}}(\ell)$ on the operational horizon.

  \item \textbf{Semantic drift.} Realistic threat trajectories from
  $\mathcal{P}_S(L)$ induce curves $\gamma_M : [0,1]\to\mathcal{M}$ with
  geodesic span bounded by a semantic drift profile $\sigma_S^{\mathbf{M}}(L)$
  as in Definition~\ref{def:semantic-drift}. In practice, this drift profile
  is estimated from historical behavior and used to distinguish benign motion
  from novel tactics.
\end{itemize}

Internally, this manifold is treated as the system's ``threat space'': genomes,
drift statistics, and risk policies all operate in $\mathcal{M}$.

\subsection{Detector and error budget mapping}

KRAKEN’s detector stack---behavior genomes, risk-aware fusion, and deterministic
replay---fits naturally into the sameness-detector abstraction of
Definition~\ref{def:sameness-detector}.

\begin{itemize}
  \item \textbf{Statistic $\Sigma_{\mathbf{M}}$.} In the geometric view, the
  detection statistic $\Sigma_{\mathbf{M}} : \mathcal{M} \to \mathbb{R}$ is
  built from (i) distances to class-specific threat behavior genomes (shrinkage
  Mahalanobis templates in $\mathcal{M}$), (ii) drift statistics along recent
  trajectories, and (iii) simple auxiliary features (environmental covariates,
  sensor trust scores). A corresponding translator-side statistic
  $\Sigma_{\mathbf{T}} : \bigsqcup_i V_i \to \mathbb{R}$ is defined by composing
  modalities through $G_S$ in the sense of Section~\ref{sec:G-construction}.

  \item \textbf{Decision rule $A$.} A CVaR-based fusion and thresholding rule
  $A$ maps the statistic and context (mission phase, alert policies) to
  $\{0,1,\mathrm{abstain}\}$: trigger, suppress, or abstain.

  \item \textbf{Error decomposition.} Following Definition~\ref{def:sameness-detector},
  KRAKEN’s end-to-end miss probability decomposes into:
  \begin{itemize}
    \item \emph{geometric/representation error} $E_{\mathrm{geom}}^{\mathbf{T}}$,
          $E_{\mathrm{geom}}^{\mathbf{M}}$ from translator drift and manifold
          distortion;
    \item \emph{linkage/observation error} $E_{\mathrm{link}}$ from noisy sensors
          and imperfect encoders;
    \item calibration error $\xi$ and abstention failure $\zeta$ from CVaR fusion
          and threshold policies;
    \item independence slack $\delta_{\mathrm{indep}}$ capturing dependence between
          failure modes.
  \end{itemize}
  Internally, these terms are estimated by a combination of held-out evaluation,
  stress tests (sensor outages, adversarial conditions), and analytic bounds of
  the type derived in Section~\ref{sec:error-transfer}.
\end{itemize}

A separate, application-focused document proves a KRAKEN-specific analogue of
the compositional error-budget theorem: under Lipschitz and margin assumptions,
stage-wise bounds on selection, translation, geometry, genomes, fusion, trust,
and replay compose into an end-to-end miss probability bound
$\varepsilon_{\mathrm{total}}$ for a given mission profile.

\subsection{How the general theorems apply}

From the perspective of this paper, the most interesting feature of KRAKEN is
that it uses \emph{both} realizations at once:
\begin{itemize}
  \item A translator-first view $\mathbf{T}_{\mathrm{mar}}$ powers cross-modal
        prediction, redundancy, and graceful degradation when some sensors fail;
  \item A manifold-first view $\mathbf{M}_{\mathrm{mar}}$ powers behavior
        genomes, drift detection, and zero-/few-shot generalization.
\end{itemize}

The functors $F_S$ and $G_S$ and the error-transfer theorems of
Section~\ref{sec:error-transfer} give a principled explanation of empirical
observations made during KRAKEN’s development:

\begin{itemize}
  \item When the translator graph passes chartability audits (small
        $\varepsilon_{\mathrm{trans}}$, $\delta_{\mathrm{asymm}}$), the induced
        manifold has bounded distortion $\varepsilon_{\mathrm{dist}}$ and
        preserves threat semantics. This is precisely the regime where
        $F_S(\mathbf{T}_{\mathrm{mar}})$ and $\mathbf{M}_{\mathrm{mar}}$ become
        $\varepsilon$-equivalent.

  \item Conversely, when geometry audits flag increased distortion or eroded
        margins in $\mathcal{M}$ (e.g., after a distribution shift), projecting
        back through $G_S$ identifies which translators, modalities, or threat
        classes are responsible. In categorical terms, this uses the stability
        of $\Delta_{\SamTrans}$ under $G_S$.

  \item Most importantly, the \emph{detection} guarantees---bounds on
        miss probability, false positive rate, and abstention behavior---are
        invariant up to first-order slack under switching between translator-
        and manifold-centric implementations, as long as chartability and
        realizability assumptions hold. This matches the observed fact that
        translator-heavy and manifold-heavy configurations of the system achieve
        comparable detection performance when tuned to the same operational
        envelope.
\end{itemize}

\subsection{Limitations and redactions}

This case study is necessarily incomplete. It deliberately omits:
\begin{itemize}
  \item detailed threat models, sensor geometries, and campaign specifics;
  \item numerical values of internal error budgets and performance metrics;
  \item implementation details of encoders, translators, genomes, and fusion
        policies.
\end{itemize}
Those belong in a separate, application-focused paper and, in some cases,
cannot be made public.

Nonetheless, even in this abstracted form, the mapping between a real, deployed
system and the semantic-sameness framework is clear:
\begin{enumerate}
  \item a well-posed sameness structure $S_{\mathrm{mar}}$;
  \item dual realizations $\mathbf{T}_{\mathrm{mar}}$ and $\mathbf{M}_{\mathrm{mar}}$
        satisfying smooth chartability and geometric realizability;
  \item an end-to-end detector whose error budget decomposes as in
        Definition~\ref{def:sameness-detector};
  \item empirical audits that instantiate the chartability and distortion checks
        required by the theorems of Section~\ref{sec:error-transfer}.
\end{enumerate}
In this sense, KRAKEN (or any similarly structured system) serves as a proof of
concept: the abstract, category-theoretic equivalence between translator-based
and manifold-based views of semantic sameness is not only mathematically sound
but also aligned with how one can design auditable, geometry-first detection
systems in high-stakes settings.

\begin{thebibliography}{99}

\bibitem{davis2025amc}
Davis, B.~R. (2025).
\emph{Adaptive Embedding Format Switching and Semantic Translation in High-Volume Telemetry Systems}.
Zenodo.
\url{https://doi.org/10.5281/zenodo.17211403}

\bibitem{davis2025herald}
Davis, B.~R. (2025).
\emph{HERALD: High-resolution Early Recognition of Antigenic Landscape Divergence}.
Zenodo.
\url{https://doi.org/10.5281/zenodo.17640400}

\bibitem{davis2025manifold}
Davis, B.~R. (2025).
\emph{The Davis Manifold: Geometry-First Detection with Compositional Error Budgets}.
Zenodo.
\url{https://doi.org/10.5281/zenodo.17642038}

\bibitem{davis2025prism}
Davis, B. (2025).
\emph{PRISM v4.1: Risk-Aware Multi-Rail Payment Reconciliation with Bounded Semantic Translation and Spectral Consensus -- Adapting NASA Safety-Critical Systems to Financial Infrastructure}.
Zenodo.
\url{https://doi.org/10.5281/zenodo.17489610}

\end{thebibliography}

\end{document}
